% concatenated from anchored_lr_geometry_v5.tex
\documentclass[11pt]{amsart}

\usepackage[margin=1in]{geometry}
\usepackage{amsmath,amssymb,amsthm,mathtools,mathrsfs}
\usepackage{enumitem}
\usepackage{hyperref}
\hypersetup{hidelinks}
\usepackage{microtype}
\usepackage{booktabs}

\usepackage{aliascnt}

\newtheorem{theorem}{Theorem}[section]

\newaliascnt{proposition}{theorem}
\newtheorem{proposition}[proposition]{Proposition}
\aliascntresetthe{proposition}

\newaliascnt{lemma}{theorem}
\newtheorem{lemma}[lemma]{Lemma}
\aliascntresetthe{lemma}

\newaliascnt{corollary}{theorem}
\newtheorem{corollary}[corollary]{Corollary}
\aliascntresetthe{corollary}

\theoremstyle{definition}

\newaliascnt{definition}{theorem}
\newtheorem{definition}[definition]{Definition}
\aliascntresetthe{definition}

\newaliascnt{convention}{theorem}
\newtheorem{convention}[convention]{Convention}
\aliascntresetthe{convention}

\newaliascnt{remark}{theorem}
\newtheorem{remark}[remark]{Remark}
\aliascntresetthe{remark}

\usepackage[capitalise,nameinlink,noabbrev]{cleveref}
\crefname{theorem}{Theorem}{Theorems}
\crefname{proposition}{Proposition}{Propositions}
\crefname{lemma}{Lemma}{Lemmas}
\crefname{corollary}{Corollary}{Corollaries}
\crefname{definition}{Definition}{Definitions}
\crefname{convention}{Convention}{Conventions}
\crefname{remark}{Remark}{Remarks}

\DeclareMathOperator{\Law}{Law}
\DeclareMathOperator{\tr}{tr}
\DeclareMathOperator{\Var}{Var}
\DeclareMathOperator{\supp}{supp}
\DeclareMathOperator{\conv}{conv}
\DeclareMathOperator{\diag}{diag}
\DeclareMathOperator{\argmax}{argmax}
\DeclareMathOperator{\Cov}{Cov}
\newcommand{\one}{\mathbf 1}
\newcommand{\R}{\mathbb R}
\newcommand{\E}{\mathbb E}
\newcommand{\PP}{\mathbb P}
\newcommand{\cP}{\mathcal P}
\newcommand{\cX}{\mathcal X}
\newcommand{\cY}{\mathcal Y}
\newcommand{\cH}{\mathcal H}
\newcommand{\cM}{\mathcal M}
\newcommand{\fM}{\mathfrak M}
\newcommand{\cK}{\mathcal K}
\newcommand{\cR}{\mathcal R}
\newcommand{\cF}{\mathcal F}
\newcommand{\cS}{\mathcal S}
\newcommand{\Id}{\mathrm{Id}}

\title[Anchored Likelihood-Ratio Geometry of Anonymous Shuffle Experiments]{Anchored Likelihood-Ratio Geometry of Anonymous Shuffle Experiments:\\Exact Privacy Envelopes and Universal Low-Budget Design}

\author{Alex Shvets}
\address{Independent Researcher, Haifa, Israel}
\email{alex@shvets.io}
\date{March 2026}

\keywords{shuffle model, differential privacy, likelihood-ratio law, affine experiment, privacy amplification, frequency estimation, minimax design, augmented randomized response}
\subjclass[2020]{62B15, 62C20, 68P27, 60E15, 62G05}

\begin{document}

\begin{abstract}
We develop a rigorous geometric framework for anonymous shuffle experiments based on a single object: an \emph{anchored affine likelihood-ratio law}, namely a mean-zero probability measure $\rho$ on the regular simplex polytope
\[
\cX_d:=\{x\in\R^{d-1}: 1+\gamma_i^\top x\ge 0\ \text{for all }i\in[d]\},
\]
where $\gamma_1,\dots,\gamma_d$ are the vertices of the centered regular simplex in $\R^{d-1}$. Every finite-output $d$-ary channel is represented, up to conditionally identical refinements, by a unique anchored law, and conversely every anchored law realizes a valid channel.

On the privacy side, pairwise shuffled privacy reduces to one-dimensional shadows of projective fibers. Fix a common local parameter $\varepsilon_0$. Among all $\varepsilon_0$-LDP channels, binary randomized response universally extremizes all convex $f$-divergences and all hockey-stick privacy profiles after shuffling, pointwise over all alphabet sizes and all ordered row pairs $i\neq j$ in the one-coordinate contamination experiment. We also prove a rigidity converse: if both directed hockey-stick envelopes are saturated at a fixed finite $n$, then the underlying pairwise one-user likelihood-ratio law is exactly the binary endpoint law.

On the design side, the covariance matrix $\Sigma(\rho)=\int xx^\top\rho(dx)$ governs the canonical estimator. Under the canonical pairwise $\chi_*$ budget, permutation averaging never increases the budget and never increases the worst-case exact i.i.d.\ or fixed-composition canonical risk. This yields exact finite-$n$ canonical optimality formulas, an exact trace-cap theorem, and a two-orbit reduction of the global $\chi_*$ frontier: every frontier point is realized by a mixture of at most two exchangeable orbit laws. In the low-budget regime, augmented randomized response is asymptotically minimax-optimal, to the sharp leading constant, over all channels and all estimators in the i.i.d.\ multinomial model.

Under the raw local $\varepsilon_0$-LDP cap, the exact anchored canonical design problem reduces to an exchangeable law within the already-known subset-selection mechanism class; the maximizing subset size is determined by an explicit one-dimensional maximization problem, and the exact finite-$n$ i.i.d.\ and fixed-composition canonical risks are explicit.

The main privacy and design arguments are self-contained and do not rely on the author's trilogy. It also clarifies a structural limitation: the anchored geometry is the correct language for exact privacy envelopes and low-budget canonical design, but it does not by itself subsume the non-Gaussian fixed-composition asymptotics of the author's trilogy.
\end{abstract}

\maketitle
\tableofcontents


\section{Introduction}

The shuffle model occupies a precise middle ground between the local and curator models of differential privacy. Each user privatizes her own datum locally, but the observer receives only an anonymized collection of the privatized messages rather than their user labels. This anonymity can substantially amplify privacy, and it also changes the geometry of optimal mechanism design. The main claim of the present paper is that, for exact privacy envelopes and exact low-budget canonical design, the right state variable is neither the raw channel matrix nor the output alphabet. It is a single mean-zero law on a fixed polytope.

\medskip
\noindent\textbf{The shuffle model.}
Fix an input alphabet $[d]=\{1,\dots,d\}$ and a finite-output local channel $W:[d]\to \Delta(\mathcal Y)$. For a dataset $x^n=(x_1,\dots,x_n)\in [d]^n$, user $m$ draws
\[
Y_m\sim W(\cdot\mid x_m)
\qquad (m=1,\dots,n)
\]
independently, and a shuffler outputs a uniformly random permutation of $(Y_1,\dots,Y_n)$. Equivalently, the observer sees only the unordered multiset, or histogram, of messages. Privacy is evaluated on neighboring datasets that differ in one coordinate. Design questions arise when the inputs are themselves random, for example i.i.d.\ with composition vector $\theta\in \Delta_d$, or deterministic with a prescribed empirical composition.

\medskip
\noindent\textbf{Privacy functionals.}
For probability measures $P,Q$ on a common space with $P\ll Q$, and a convex function $f:[0,\infty)\to \R$ with $f(1)=0$, write
\[
D_f(P\|Q):=\int f\!\left(\frac{dP}{dQ}\right)dQ.
\]
We also use the hockey-stick functional
\[
H_\alpha(P,Q):=\int\left(\frac{dP}{dQ}-\alpha\right)_+ dQ,\qquad \alpha\ge 1.
\]
The exact one-sided $(\varepsilon,\delta)$ tradeoff is encoded by $H_{e^\varepsilon}(P,Q)$, and the usual two-sided differential privacy condition requires the same bound after exchanging $P$ and $Q$. Thus a pointwise upper envelope for all hockey-stick profiles is stronger than a single asymptotic $(\varepsilon,\delta)$ estimate. One purpose of the present paper is to obtain such envelopes exactly at finite $n$, uniformly over all alphabet sizes and all local mechanisms satisfying a prescribed local constraint.

\medskip
\noindent\textbf{The anchored law.}
Let $H\in \R^{d\times(d-1)}$ have orthonormal columns spanning
\[
T_d:=\{u\in \R^d:\one^\top u=0\},
\]
and define
\[
\gamma_i:=H^\top e_i\in \R^{d-1},\qquad i=1,\dots,d.
\]
Then $\gamma_1,\dots,\gamma_d$ are the vertices of the centered regular simplex:
\[
\sum_{i=1}^d \gamma_i=0,
\qquad
\gamma_i^\top \gamma_j=\delta_{ij}-\frac1d.
\]
The corresponding simplex polytope is
\[
\cX_d:=\{x\in \R^{d-1}:1+\gamma_i^\top x\ge 0\ \text{for all }i\in[d]\}.
\]
An \emph{anchored affine likelihood-ratio law} is a probability measure $\rho\in \cP(\cX_d)$ with mean zero. The exact representation theorem in Section~\ref{sec:anchored} shows that every finite-output $d$-ary channel is represented by a unique such law, up to conditionally identical refinements. In this gauge, one-user mixture laws become affine tilts of $\rho$, pairwise likelihood ratios become scalar shadows of projective fibers, and canonical design becomes a moment problem for the covariance matrix
\[
\Sigma(\rho):=\int xx^\top\,\rho(dx).
\]
The representation itself is a coordinate-specific instance of the classical posterior/standard-experiment representation of statistical experiments; see Blackwell \cite{blackwell-comparison} and Torgersen \cite[Chapter~2]{torgersen}. The contribution is the explicit regular-simplex gauge adapted to the shuffle privacy and canonical design problems below.

\medskip
\noindent\textbf{Why this paper is not the trilogy in disguise.}
The author's trilogy \cite{shvets-part1,shvets-part2,shvets-part3} studies neighboring shuffle experiments through exact likelihood-ratio identities, quotient carriers, and Gaussian/non-Gaussian asymptotic regimes. The present paper has a different scope. It works with a single $d$-ary channel rather than a neighboring two-row pair, and it treats exact finite-$n$ privacy envelopes and exact canonical design problems. The anchored law is the correct carrier for those exact questions, but it is not a replacement for the quotient geometry that drives the non-Gaussian fixed-composition asymptotics of Parts~II--III. Making that division of labor mathematically explicit is one of the goals of this paper.

A second point of comparison is the paper \cite{shvets-growing} on growing alphabets and low-budget design. That work established several headline phenomena: universal binary extremality at the $\chi^2$ level, asymptotic low-budget optimality of augmented randomized response under the $\chi_*$-budget, optimality of GRR within the subset-selection family, and a van Trees lower bound. More concretely, \cite{shvets-growing} already proved exact pairwise likelihood-ratio compression, the sharp universal pairwise $\chi^2$ extremal bound with endpoint two-point equality characterization, and low-budget augmented-GRR optimality among permutation-equivariant finite-output channels; what is new here is the anchored-law reformulation, the finite-$n$ convex-order/hockey-stick upgrade with rigidity, the sharp all-estimators low-budget constant, and the closure of the full global $\chi_*$ frontier.

\medskip
\noindent\textbf{Related work.}
The shuffle-model privacy literature has developed several complementary analytic frameworks. The foundational shuffle-model amplification results were established by Erlingsson, Feldman, Mironov, Raghunathan, Talwar, and Thakurta \cite{erlingsson-shuffle} and, independently, by Cheu, Smith, Ullman, Zeber, and Zhilyaev \cite{cheu-shuffle}. The privacy-blanket decomposition of Balle, Bell, Gasc\'on, and Nissim \cite{balle-blanket} provides a general method for amplification bounds. Feldman, McMillan, and Talwar introduced the clone paradigm and then sharpened it to obtain stronger approximate-DP and R\'enyi-DP amplification bounds together with improved numerical accounting \cite{feldman-clones,feldman-stronger}. Girgis, Data, Diggavi, Suresh, and Kairouz derived R\'enyi-DP guarantees for general discrete local randomizers in shuffled distributed learning \cite{girgis-rdp}. Koskela, Heikkil\"a, and Honkela \cite{koskela-shuffle} developed numerical accounting methods for the shuffle model based on hockey-stick divergences. Wang et al.\ \cite{wang-unified} gave a unified variation-ratio reduction framework with tight numerical amplification bounds. More recently, Takagi and Liew \cite{takagi-liew} analyzed shuffling beyond pure local differential privacy through blanket divergence, introduced a one-parameter shuffle index, and obtained asymptotic upper and lower bands together with an FFT-based finite-$n$ accountant.

On the mechanism-design side, Kairouz, Oh, and Viswanath \cite{kov} showed that for a broad class of convex utility functionals under pure local differential privacy one may restrict attention to staircase mechanisms; they also identified binary/randomized-response-type mechanisms as universally optimal in the appropriate privacy regimes. For frequency estimation under local differential privacy, Wang, Blocki, Li, and Jha \cite{wang-ldp} introduced subset-selection mechanisms, and Ye and Barg \cite{ye-barg} proved that subset selection with an optimized subset size is asymptotically order-optimal.

The present paper is complementary to all of the above. Its central objects are an anchored affine law, projective fibers, and convex-order arguments. This leads to exact finite-$n$ universal envelopes for all convex $f$-divergences and all hockey-stick profiles, together with a rigidity converse, and to exact canonical mechanism design statements under two natural local budgets. To our knowledge, neither the projective-fiber representation nor the converse rigidity statement appears in the blanket, clone, or extremal-mechanism formalisms in this exact form. On the raw-cap design side, the optimizer class---subset selection---was identified in earlier work \cite{wang-ldp,ye-barg}; the present contribution is the exact anchored canonical formulation, explicit finite-$n$ risk formulas, and the integration with the shuffle-privacy envelope theory.

\medskip
\noindent\textbf{Main contributions.}
The paper makes five theorem-level contributions.

\begin{enumerate}[label=\textup{(\arabic*)}]
\item \emph{Anchored representation as a gauge.}
We specialize the classical posterior/standard-experiment representation to the regular-simplex gauge adapted to shuffle privacy and canonical design. Every finite-output $d$-ary channel induces a unique finitely supported anchored law $\rho\in\fM_d$ up to conditionally identical refinements, and every anchored law defines a valid $\cX_d$-valued channel.

\item \emph{Projective fiber reduction for privacy.}
Pairwise shuffled privacy reduces exactly to one-dimensional shadows of projective fibers. The core one-dimensional pairwise reduction was established in \cite{shvets-growing}; the present formulation embeds it in the global projective geometry of Section~\ref{sec:projective}, which additionally provides the transport formula (\cref{prop:fiber-props}(ii)), the exact interior converse (\cref{thm:converse-fiber}), and the boundary analysis (\cref{rem:boundary-fibers}).

\item \emph{Universal convex-order privacy envelope and rigidity.}
Building on the pairwise likelihood-ratio reduction, we show that the binary endpoint law maximizes after shuffling all convex functionals of the neighboring likelihood-ratio average under a common $\varepsilon_0$-LDP cap, hence all convex $f$-divergences and both directed hockey-stick profiles. We also prove a finite-$n$ rigidity converse: simultaneous saturation of both directed hockey-stick envelopes forces the one-user pairwise likelihood-ratio law to be the endpoint two-point law.

\item \emph{Exact canonical design under the $\chi_*$ budget.}
We derive exact i.i.d.\ and fixed-composition canonical risk formulas, prove the trace cap, and show that the full global $\chi_*$ frontier is the upper concave envelope of the orbit curve $\Gamma_d$; consequently every frontier point is realized by a mixture of at most two exchangeable orbit laws. In the low-budget regime this recovers augmented randomized response as the unique canonical optimizer.

\item \emph{Exact canonical design under the raw local cap.}
Under the raw local $\varepsilon_0$-LDP constraint, the exact canonical design problem admits an exchangeable subset-selection solution---a mechanism class identified in earlier work \cite{wang-ldp,ye-barg}. Every extreme exchangeable optimizer is a maximizing subset-selection law, and every exchangeable optimizer is a convex combination of the maximizing subset-selection laws. The present paper gives the exact anchored canonical formulation: the optimizing subset size is determined by the one-dimensional family $s\mapsto T_{d,e^{\varepsilon_0}}(s)$, ties are handled explicitly, and the exact finite-$n$ canonical risks are closed-form.
\end{enumerate}

\medskip
\noindent\textbf{A notation bridge to the trilogy.}
Readers familiar with \cite{shvets-part1,shvets-part2,shvets-part3} may find it useful to keep the following dictionary in mind. In the trilogy, the basic object is a neighboring two-row experiment and its quotient likelihood field. In the present paper, the single anchored law $\rho$ plays the role of the exact carrier for a full $d$-ary channel. The centered composition coordinate is
\[
h_\theta=H^\top\!\left(\theta-\frac{\one}{d}\right)\in \cH_d,
\]
which is the anchored image of the composition vector. The covariance matrix $\Sigma(\rho)$ is an exact canonical-design object; it should not be confused with the composition-dependent neighboring covariance matrices $\Sigma_\pi$ of the trilogy. Likewise, the projective fibers $\eta_h$ are exact finite-$n$ objects attached to the anchored law, not the non-Gaussian quotient carriers of Parts~II--III.

\medskip
\noindent\textbf{Scope and limitation.}
The present paper is self-contained on the privacy and design sides. No result from the trilogy is used in any proof. At the same time, the scope is deliberately narrower than a full asymptotic unification theorem. For interior base points, projective fibers are bounded, so any asymptotic theory based only on the interior anchored gauge is necessarily Gaussian there. That observation is structural, not cosmetic: it explains why the non-Gaussian fixed-composition limits of \cite{shvets-part2,shvets-part3} require a different carrier.

\medskip
\noindent\textbf{Organization of the paper.}
Section~\ref{sec:anchored} establishes the exact anchored representation. Section~\ref{sec:projective} develops projective maps and fibers, with an explicit separation between interior and boundary base points. Section~\ref{sec:privacy} proves the one-dimensional privacy reduction, the universal binary envelope, and the rigidity converse. Sections~\ref{sec:design}--\ref{sec:tracecap} develop the canonical estimator, global symmetrization, the exact trace cap, and the two-orbit frontier reduction. Section~\ref{sec:asymptotic} proves universal low-budget asymptotic minimax optimality. Section~\ref{sec:raw} solves the exact canonical design problem under the raw local cap and records a numerical illustration. Section~\ref{sec:discussion} discusses scope and open problems.

\medskip
\noindent\textbf{Notation.}
We write $[d]=\{1,\dots,d\}$, $\Delta_d=\{\theta\in \R_+^d:\sum_i \theta_i=1\}$, and
\[
\cH_d:=\conv\{\gamma_1,\dots,\gamma_d\}\subset \R^{d-1}.
\]
When no confusion is possible we write $\cH=\cH_d$.


\section{Anchored affine likelihood-ratio laws}\label{sec:anchored}

\subsection{Geometry of the simplex polytope}

\begin{definition}[Regular simplex data]
Fix $H\in\R^{d\times(d-1)}$ with orthonormal columns spanning $T_d=\{u\in\R^d:\one^\top u=0\}$. Define
\[
\gamma_i:=H^\top e_i,\qquad i=1,\dots,d,
\]
and
\[
\cX_d:=\{x\in\R^{d-1}:1+\gamma_i^\top x\ge 0\text{ for all }i\in[d]\}.
\]
Define also the template simplex
\[
\cK_d:=\{a\in\R_+^d:\sum_{i=1}^d a_i=d\}.
\]
\end{definition}

\begin{lemma}[Affine identification of $\cX_d$ and $\cK_d$]\label{lem:polytope-template}
The map
\[
a(x):=\one+Hx\in\R^d
\]
is an affine bijection from $\cX_d$ onto $\cK_d$, with inverse
\[
x(a)=H^\top(a-\one).
\]
Moreover,
\[
a_i(x)=1+\gamma_i^\top x,
\qquad
\|x\|_2^2=\|a(x)-\one\|_2^2=\sum_{i=1}^d a_i(x)^2-d.
\]
In particular,
\[
\sup_{x\in\cX_d}\|x\|_2=\sqrt{d(d-1)}.
\]
\end{lemma}

\begin{proof}
Since $H^\top\one=0$ and $H^\top H=I_{d-1}$, we have
\[
\sum_{i=1}^d a_i(x)=\one^\top(\one+Hx)=d.
\]
Also
\[
a_i(x)=e_i^\top(\one+Hx)=1+e_i^\top Hx=1+\gamma_i^\top x,
\]
so $x\in\cX_d$ iff $a(x)\in\cK_d$. Conversely, if $a\in\cK_d$, then $a-\one\in T_d$ and therefore
\[
H H^\top(a-\one)=a-\one,
\]
which yields
\[
a=\one+H H^\top(a-\one)=a(x(a)).
\]
Thus $x\mapsto a(x)$ is a bijection with inverse $a\mapsto H^\top(a-\one)$.

Finally,
\[
\|x\|_2^2=\|H^\top(a-\one)\|_2^2=\|a-\one\|_2^2
\]
because $a-\one\in T_d$ and $H^\top$ is an isometry on $T_d$. Since
\[
\|a-\one\|_2^2=\sum_{i=1}^d(a_i^2-2a_i+1)=\sum_{i=1}^d a_i^2-d,
\]
the identity follows. The maximum over $\cK_d$ is attained at a vertex $de_i$, giving
\[
\max_{a\in\cK_d}\|a-\one\|_2^2=(d-1)^2+(d-1)=d(d-1).
\]
\end{proof}

\begin{definition}[Anchored law]
An \emph{anchored affine likelihood-ratio law} is a probability measure $\rho\in\cP(\cX_d)$ such that
\[
\int_{\cX_d} x\,\rho(dx)=0.
\]
We write
\[
\fM_d:=\Bigl\{\rho\in\cP(\cX_d):\int x\,\rho(dx)=0\Bigr\}.
\]
\end{definition}

\subsection{Exact representation and reconstruction}

\begin{definition}[Conditionally identical refinements]\label{def:cond-id-ref}
Two channels $W:[d]\to\Delta(\cY)$ and $W':[d]\to\Delta(\cY')$ are said to differ only by \emph{conditionally identical refinements} if one can be obtained from the other by splitting or merging output symbols in a way that does not depend on the input row: concretely, there exists a channel $V:\cY\to\Delta(\cY')$ such that $W'(B\mid i)=\sum_y W(y\mid i)\,V(B\mid y)$ for all $i$ and all measurable $B\subseteq\cY'$, and conversely. Equivalently, $W$ and $W'$ induce the same statistical experiment as viewed through Le~Cam sufficiency.
\end{definition}

\begin{theorem}[Exact affine representation]\label{thm:representation}
Let $W:[d]\to\Delta(\mathcal Y)$ be a finite-output channel.

\begin{enumerate}[label=\textup{(\roman*)}]
\item Define the row average
\[
\bar W(y):=\frac1d\sum_{i=1}^d W(y\mid i),
\]
and, whenever $\bar W(y)>0$,
\[
a_i(y):=\frac{W(y\mid i)}{\bar W(y)},\qquad x(y):=H^\top(a(y)-\one).
\]
Then $x(y)\in\cX_d$ and the pushforward law
\[
\rho_W:=\sum_{y:\bar W(y)>0}\bar W(y)\,\delta_{x(y)}
\]
belongs to $\fM_d$.

\item Let $W_{\rho_W}$ be the channel on $\cX_d$ defined by
\[
W_{\rho_W}(B\mid i):=\int_B (1+\gamma_i^\top x)\,\rho_W(dx).
\]
Then $W_{\rho_W}$ is the exact sufficient reduction of $W$ obtained by merging output symbols with the same anchored coordinate. In particular, $W$ and $W_{\rho_W}$ differ only by conditionally identical refinements.

\item Conversely, every $\rho\in\fM_d$ defines a valid channel
\[
W_\rho(B\mid i):=\int_B (1+\gamma_i^\top x)\,\rho(dx),\qquad B\subseteq\cX_d,
\]
and the map $W\mapsto\rho_W$ descends to a bijection between finite-output $d$-ary channels modulo conditionally identical refinements and the finitely supported members of $\fM_d$. More generally, every $\rho\in\fM_d$ defines a valid $\cX_d$-valued channel $W_\rho$.

\item For any composition $\theta\in\Delta_d$, define
\[
h_\theta:=H^\top\Bigl(\theta-\frac{\one}{d}\Bigr)=\sum_{i=1}^d \theta_i\gamma_i\in\cH.
\]
Then the one-user mixture law under input composition $\theta$ is
\[
q_\theta(dx)=\bigl(1+h_\theta^\top x\bigr)\rho(dx).
\]
\end{enumerate}
\end{theorem}

\begin{proof}
(i) Since $\sum_i a_i(y)=d$, \cref{lem:polytope-template} implies $x(y)\in\cX_d$. Also
\[
\int x\,\rho_W(dx)=\sum_y \bar W(y) H^\top(a(y)-\one)
=H^\top\Bigl(\sum_y \bar W(y)a(y)-\one\Bigr).
\]
For every coordinate $i$,
\[
\sum_y \bar W(y) a_i(y)=\sum_y W(y\mid i)=1,
\]
so $\sum_y \bar W(y)a(y)=\one$, and therefore $\int x\,\rho_W(dx)=0$.

(ii) Let $B\subseteq\cX_d$ be Borel. Then
\[
\begin{aligned}
W_{\rho_W}(B\mid i)
&=\sum_{y:\bar W(y)>0}\mathbf 1_{\{x(y)\in B\}}\,(1+\gamma_i^\top x(y))\,\bar W(y)\\
&=\sum_{y:\bar W(y)>0}\mathbf 1_{\{x(y)\in B\}}\,a_i(y)\,\bar W(y)\\
&=\sum_{y:\bar W(y)>0}\mathbf 1_{\{x(y)\in B\}}\,W(y\mid i).
\end{aligned}
\]
Thus $W_{\rho_W}$ is obtained by merging all output symbols with identical anchored coordinate $x(y)$; within each such fiber the conditional law given the reduced output is independent of the input row. That is exactly a conditionally identical refinement.

(iii) For any $\rho\in\fM_d$, nonnegativity follows from $1+\gamma_i^\top x\ge 0$ on $\cX_d$, and
\[
W_\rho(\cX_d\mid i)=\int (1+\gamma_i^\top x)\,\rho(dx)=1+\gamma_i^\top\int x\,\rho(dx)=1.
\]
So $W_\rho$ is a channel. Moreover, the row average recovers $\rho$ exactly:
\[
\frac1d\sum_{i=1}^d W_\rho(dx\mid i)=\frac1d\sum_{i=1}^d (1+\gamma_i^\top x)\rho(dx)=\rho(dx)
\]
because $\sum_i \gamma_i=0$. Equivalently,
\[
\rho=\frac1d\sum_{i=1}^d W_\rho(\cdot\mid i).
\]
If $\rho$ is finitely supported, then $W_\rho$ is a finite-output channel on the support of $\rho$. For any support point $x$,
\[
\bar W_\rho(\{x\})=\frac1d\sum_{i=1}^d W_\rho(\{x\}\mid i)=\rho(\{x\}),
\]
and therefore
\[
a_i(x)=\frac{W_\rho(\{x\}\mid i)}{\bar W_\rho(\{x\})}=1+\gamma_i^\top x.
\]
Applying the construction from part~(i) to $W_\rho$ thus returns the same anchored coordinate $x$, so $\rho_{W_\rho}=\rho$. Together with part~(ii), this proves the claimed bijection on equivalence classes. Uniqueness of $\rho$ for a given $W_\rho$ is exactly the row-average identity above.

(iv) This is immediate:
\[
\sum_{i=1}^d \theta_i W_\rho(dx\mid i)
=\sum_i \theta_i(1+\gamma_i^\top x)\rho(dx)
=\bigl(1+h_\theta^\top x\bigr)\rho(dx).
\qedhere
\]
\end{proof}


\section{Projective transport and fibers}\label{sec:projective}

\subsection{Projective maps}

\begin{definition}[Projective domains and maps]
For $h\in\cH$, define
\[
A_h:=\{x\in\cX_d:1+h^\top x>0\},
\qquad
\cY_h:=\{y\in\R^{d-1}:1+(\gamma_i-h)^\top y\ge 0\ \forall i\in[d]\},
\]
and
\[
T_h(x):=\frac{x}{1+h^\top x},\qquad x\in A_h,
\]
\[
S_h(y):=\frac{y}{1-h^\top y},\qquad y\in\cY_h.
\]
\end{definition}

\begin{proposition}[Projective bijection]\label{prop:diffeo}
For every $h\in\cH$, $T_h:A_h\to\cY_h$ is a bijection with inverse $S_h$. Moreover:
\begin{enumerate}[label=\textup{(\roman*)}]
\item $1-h^\top y>0$ for every $y\in\cY_h$;
\item $S_h(\cY_h)\subseteq\cX_d$ and $T_h(A_h)\subseteq\cY_h$;
\item if $h\in\operatorname{int}\cH$, then $A_h=\cX_d$;
\item if $h\in\partial \cH$, then $\cY_h$ is unbounded.
\end{enumerate}
\end{proposition}

\begin{proof}
Let $x\in A_h$. Then for each $i$,
\[
1+(\gamma_i-h)^\top T_h(x)
=1+(\gamma_i-h)^\top \frac{x}{1+h^\top x}
=\frac{1+\gamma_i^\top x}{1+h^\top x}\ge 0.
\]
Hence $T_h(x)\in\cY_h$.

Now let $y\in\cY_h$. Summing the inequalities $1+(\gamma_i-h)^\top y\ge 0$ over $i$ gives
\[
d+\Bigl(\sum_{i=1}^d \gamma_i-dh\Bigr)^\top y=d(1-h^\top y)\ge 0.
\]
If equality held, then each summand would be zero, so
\[
1+(\gamma_i-h)^\top y=0\qquad \forall i.
\]
Subtracting two such equations yields $(\gamma_i-\gamma_j)^\top y=0$ for all $i,j$. Since the differences $\gamma_i-\gamma_j$ span $\R^{d-1}$, this implies $y=0$, impossible. Thus $1-h^\top y>0$.

Using that positivity,
\[
1+\gamma_i^\top S_h(y)
=1+\gamma_i^\top \frac{y}{1-h^\top y}
=\frac{1+(\gamma_i-h)^\top y}{1-h^\top y}\ge 0,
\]
so $S_h(y)\in\cX_d$. Furthermore,
\[
1+h^\top S_h(y)=1+\frac{h^\top y}{1-h^\top y}=\frac{1}{1-h^\top y}>0,
\]
so $S_h(y)\in A_h$.

A direct calculation gives
\[
T_h(S_h(y))=\frac{y/(1-h^\top y)}{1+h^\top y/(1-h^\top y)}=y,
\]
and similarly $S_h(T_h(x))=x$. Thus $T_h$ and $S_h$ are inverse bijections.

If $h=\sum_i \theta_i\gamma_i$ with all $\theta_i>0$, then for $x\in\cX_d$,
\[
1+h^\top x=\sum_{i=1}^d \theta_i(1+\gamma_i^\top x).
\]
The summands are nonnegative and cannot all vanish simultaneously, so the sum is strictly positive. Hence $A_h=\cX_d$.

Finally, if $h\in\partial \cH$, choose a supporting hyperplane of $\cH$ at $h$ with normal vector $u\neq 0$, so that $u^\top z\le u^\top h$ for all $z\in \cH$. Setting $r:=-u$, we obtain
\[
(\gamma_i-h)^\top r\ge 0\qquad (i=1,\dots,d).
\]
Equivalently, $r$ is a recession direction of $\cY_h$. Hence for every $t\ge 0$,
\[
1+(\gamma_i-h)^\top (tr)=1+t(\gamma_i-h)^\top r\ge 0,
\]
so $tr\in \cY_h$. Thus $\cY_h$ is unbounded.
\end{proof}

\subsection{Projective fibers and transport}

\begin{definition}[Fibers]
Let $\rho\in\fM_d$ and $h\in\cH$. Define
\[
q_h(dx):=(1+h^\top x)\rho(dx).
\]
Since $q_h(\cX_d)=1$, this is a probability measure. Its projective fiber is
\[
\eta_h:=(T_h)_\# q_h\in\cP(\cY_h).
\]
\end{definition}

\begin{proposition}[Interior fiber properties and transport]\label{prop:fiber-props}
Let $\rho\in\fM_d$, let $h\in\operatorname{int}\cH$, and let $k\in\cH$.
\begin{enumerate}[label=\textup{(\roman*)}]
\item $\eta_h$ has mean zero:
\[
\int y\,\eta_h(dy)=0.
\]
\item The fibers are related by the exact transport formula
\[
\eta_k=(T_{k-h})_\#\bigl((1+(k-h)^\top y)\eta_h(dy)\bigr),
\]
where $T_{k-h}(y):=y/(1+(k-h)^\top y)$.
\item $\eta_h$ has bounded support.
\end{enumerate}
\end{proposition}

\begin{proof}
Because $h\in \operatorname{int}\cH$, \cref{prop:diffeo}(iii) gives $A_h=\cX_d$.

(i) By definition and \cref{prop:diffeo},
\[
\int y\,\eta_h(dy)
=
\int T_h(x)\,q_h(dx)
=
\int \frac{x}{1+h^\top x}(1+h^\top x)\,\rho(dx)
=
\int x\,\rho(dx)=0.
\]

(ii) Since $q_k$ vanishes on $F_k:=\{x:1+k^\top x=0\}$, we may integrate over $A_k$. For $x\in A_k$ one has
\[
T_k(x)=\frac{x}{1+k^\top x}
=
\frac{x/(1+h^\top x)}{1+(k-h)^\top x/(1+h^\top x)}
=
T_{k-h}(T_h(x)),
\]
and
\[
1+(k-h)^\top T_h(x)=\frac{1+k^\top x}{1+h^\top x}.
\]
Therefore, for any bounded measurable $\varphi$,
\[
\begin{aligned}
\int \varphi(z)\,\eta_k(dz)
&=
\int_{A_k}\varphi(T_k(x))(1+k^\top x)\,\rho(dx)\\
&=
\int_{A_k}\varphi(T_{k-h}(T_h(x)))\bigl(1+(k-h)^\top T_h(x)\bigr)(1+h^\top x)\,\rho(dx)\\
&=
\int \varphi(T_{k-h}(y))\bigl(1+(k-h)^\top y\bigr)\,\eta_h(dy).
\end{aligned}
\]
This is the stated formula. (When $k\in\partial\cH$, the map $T_{k-h}$ is applied only on the set $\{y:1+(k-h)^\top y>0\}$; on its complement the weight $1+(k-h)^\top y$ vanishes, so the weighted measure $(1+(k-h)^\top y)\,\eta_h(dy)$ is zero there regardless of whether $\eta_h$ charges that set.)

(iii) Since $\cX_d$ is compact and $1+h^\top x$ is continuous and strictly positive on $\cX_d$, there exists $c_h>0$ with $1+h^\top x\ge c_h$ on $\cX_d$. Hence
\[
\|T_h(x)\|\le c_h^{-1}\|x\|\le c_h^{-1}\sqrt{d(d-1)},
\]
so $T_h(\cX_d)$ is bounded. Since $\eta_h$ is supported on $T_h(\cX_d)$, it has bounded support.
\end{proof}

\begin{remark}[Boundary fibers can lose mass]\label{rem:boundary-fibers}
When $h\in\partial \cH$, the map $T_h$ is defined only on $A_h=\cX_d\setminus F_h$, where
\[
F_h:=\{x\in \cX_d:1+h^\top x=0\}.
\]
Any mass of $\rho$ carried by $F_h$ is invisible to
\[
q_h(dx)=(1+h^\top x)\rho(dx)
\]
and therefore invisible to the fiber $\eta_h=(T_h)_\# q_h$. Consequently $\eta_h$ need not have mean zero, and the fiber does not determine $\rho$ uniquely unless one imposes the additional condition $\rho(F_h)=0$.
\end{remark}

\begin{theorem}[Exact converse realizability for interior fibers]\label{thm:converse-fiber}
Fix $h\in\operatorname{int}\cH$. Let $\eta\in \cP(\cY_h)$ have finite first moment (automatically satisfied when $h\in\operatorname{int}\cH$, since $\cY_h=T_h(\cX_d)$ is bounded as shown in the proof of \cref{prop:fiber-props}(iii)) and satisfy
\[
\int y\,\eta(dy)=0.
\]
Define a measure on $\cX_d$ by
\[
\rho:=(S_h)_\#\bigl((1-h^\top y)\eta(dy)\bigr).
\]
Then $\rho\in \fM_d$, it is the unique anchored law with fiber $\eta_h=\eta$, and for its base-point law
\[
q_h(dx)=(1+h^\top x)\rho(dx)
\]
one has
\[
(T_h)_\# q_h=\eta.
\]
\end{theorem}

\begin{proof}
By \cref{prop:diffeo}, $1-h^\top y>0$ on $\cY_h$, so $\rho$ is a positive measure. Its total mass is
\[
\rho(\cX_d)=\int (1-h^\top y)\,\eta(dy)=1-h^\top \int y\,\eta(dy)=1.
\]
Thus $\rho$ is a probability measure. Since $h\in \operatorname{int}\cH$, one also has
\[
\rho(A_h)=\int (1-h^\top y)\,\eta(dy)=1,
\]
because $A_h=\cX_d$.

Since $S_h(\cY_h)\subseteq \cX_d$, the support is correct. Its mean is
\[
\int x\,\rho(dx)=\int S_h(y)(1-h^\top y)\,\eta(dy)=\int y\,\eta(dy)=0,
\]
so $\rho\in \fM_d$.

Now let $\varphi$ be bounded and measurable on $\cY_h$. Then
\[
\begin{aligned}
\int \varphi(T_h(x))\,q_h(dx)
&=
\int \varphi(T_h(x))(1+h^\top x)\,\rho(dx)\\
&=
\int \varphi(T_h(S_h(y)))\,(1+h^\top S_h(y))(1-h^\top y)\,\eta(dy).
\end{aligned}
\]
But $T_h(S_h(y))=y$ and
\[
1+h^\top S_h(y)=1+\frac{h^\top y}{1-h^\top y}=\frac{1}{1-h^\top y}.
\]
Therefore
\[
\int \varphi(T_h(x))\,q_h(dx)=\int \varphi(y)\,\eta(dy),
\]
which proves $(T_h)_\# q_h=\eta$.

For uniqueness, suppose $\rho'$ is another anchored law with the same fiber. Since $h\in \operatorname{int}\cH$, the map $T_h:\cX_d\to \cY_h$ is bijective. Therefore
\[
(T_h)_\# q_h'=(T_h)_\# q_h=\eta
\qquad\Longrightarrow\qquad
q_h'=q_h,
\]
where $q_h'(dx)=(1+h^\top x)\rho'(dx)$. Because $1+h^\top x>0$ on all of $\cX_d$, division by $1+h^\top x$ yields
\[
\rho'(dx)=\frac{q_h'(dx)}{1+h^\top x}=\frac{q_h(dx)}{1+h^\top x}=\rho(dx).
\]
Thus $\rho=\rho'$.
\end{proof}


\section{Pairwise privacy reduction and universal envelopes}\label{sec:privacy}

\subsection{Exact one-dimensional reduction}

\begin{proposition}[Pairwise likelihood ratios are scalar shadows]\label{prop:pairwise-lr}
Let $h,k\in\cH$, let $\rho\in\fM_d$, and assume $q_k\ll q_h$. Let $Y_h\sim\eta_h$. Then
\[
\frac{dq_k}{dq_h}(x)=\frac{1+k^\top x}{1+h^\top x}=1+(k-h)^\top T_h(x)
\qquad q_h\text{-a.s.}
\]
Hence, under $q_h$, the one-user likelihood ratio between $q_k$ and $q_h$ has law
\[
\mu_{h,k}:=\Law\bigl(1+(k-h)^\top Y_h\bigr).
\]
In particular, if $W_\rho$ is $\varepsilon_0$-LDP (so that all row pairs are mutually absolutely continuous), then for a row pair $i\neq j$,
\[
L_{ij}:=\frac{dW_\rho(\cdot\mid j)}{dW_\rho(\cdot\mid i)}(X)=1+(\gamma_j-\gamma_i)^\top Y_{\gamma_i},\qquad X\sim W_\rho(\cdot\mid i).
\]
If
\[
Q_{0,n}^{ij}:=W_\rho(\cdot\mid i)^{\otimes n},
\qquad
Q_{1,n}^{ij}:=\frac1n\sum_{m=1}^n W_\rho(\cdot\mid i)^{\otimes(m-1)}\otimes W_\rho(\cdot\mid j)\otimes W_\rho(\cdot\mid i)^{\otimes(n-m)},
\]
then under $Q_{0,n}^{ij}$,
\[
\frac{dQ_{1,n}^{ij}}{dQ_{0,n}^{ij}}=\frac1n\sum_{m=1}^n L_{ij,m},
\]
where $L_{ij,1},\dots,L_{ij,n}$ are i.i.d.\ with law $\mu_{ij}:=\Law(L_{ij})$.
\end{proposition}

\begin{proof}
The first identity is immediate:
\[
\frac{1+k^\top x}{1+h^\top x}=1+\frac{(k-h)^\top x}{1+h^\top x}=1+(k-h)^\top T_h(x).
\]
Thus the one-user likelihood ratio under $q_h$ is a scalar projection of $Y_h$.

For the shuffled neighboring experiment, write
\[
L_{ij}(x):=\frac{dW_\rho(\cdot\mid j)}{dW_\rho(\cdot\mid i)}(x).
\]
Then
\[
\frac{dQ_{1,n}^{ij}}{dQ_{0,n}^{ij}}(x_1,\dots,x_n)
=\frac1n\sum_{m=1}^n L_{ij}(x_m),
\]
which is the standard likelihood-ratio formula for a one-coordinate contamination mixture. The right-hand side is symmetric in $(x_1,\dots,x_n)$, so it is measurable with respect to the unordered sample as well; hence the same likelihood-ratio expression governs the shuffled experiment. Under $Q_{0,n}^{ij}$ the coordinates are i.i.d.\ from $W_\rho(\cdot\mid i)$, so the summands are i.i.d.\ with law $\mu_{ij}$.
\end{proof}


\begin{remark}[Boundary base points in the privacy reduction]\label{rem:privacy-boundary}
For the row pair $(i,j)$ in \cref{prop:pairwise-lr}, the base point is $h=\gamma_i\in \partial\cH$. The proposition uses only the definition of the fiber as the pushforward of $q_h$ and the projective identity from \cref{prop:diffeo}; it does \emph{not} use the mean-zero statement or the converse theorem from Section~\ref{sec:projective}. Thus the universal envelope and rigidity results below are unaffected by the boundary pathology described in \cref{rem:boundary-fibers}.
\end{remark}

\begin{remark}[Relation to the quotient compression of {\cite{shvets-growing}}]\label{rem:growing-relation}
The one-dimensional pairwise reduction itself---that the shuffled neighboring experiment depends only on the pairwise likelihood-ratio law---was established in a different form in \cite[Theorem~3.2]{shvets-growing}. The present formulation embeds this fact in the global projective geometry of Section~\ref{sec:projective}, which additionally provides the transport formula (\cref{prop:fiber-props}(ii)), the exact interior converse (\cref{thm:converse-fiber}), and the boundary analysis (\cref{rem:boundary-fibers}).
\end{remark}

\begin{remark}[Absolute continuity under a common local cap]
Under a common $\varepsilon_0$-LDP cap, all row pairs are mutually absolutely continuous, so \cref{prop:pairwise-lr} applies without further qualification in \cref{thm:privacy-envelope,thm:privacy-rigidity}.
\end{remark}



\subsection{The universal binary envelope}


\begin{definition}[The bounded mean-one class]
For $\varepsilon_0\ge 0$, set
\[
a:=e^{-\varepsilon_0},\qquad b:=e^{\varepsilon_0},
\]
and define
\[
\cM_{\varepsilon_0}:=\bigl\{\mu\in\cP([a,b]): \int z\,\mu(dz)=1\bigr\}.
\]
If $\varepsilon_0=0$, define
\[
\mu^\star:=\delta_1.
\]
If $\varepsilon_0>0$, let $\mu^\star\in\cM_{\varepsilon_0}$ be the unique two-point law on $\{a,b\}$ with mean $1$:
\[
\mu^\star=\frac{b-1}{b-a}\,\delta_a+\frac{1-a}{b-a}\,\delta_b.
\]
\end{definition}


\begin{lemma}[Convex-order maximum on $\cM_{\varepsilon_0}$]\label{lem:cx-max}
For every $\mu\in\cM_{\varepsilon_0}$ and every convex function $\phi:[a,b]\to\R$,
\[
\int \phi(z)\,\mu(dz)\le \int \phi(z)\,\mu^\star(dz).
\]
Equivalently, $\mu\preceq_{cx}\mu^\star$.
\end{lemma}

\begin{proof}
If $\varepsilon_0=0$, then $a=b=1$, $\cM_{\varepsilon_0}=\{\delta_1\}$, and the claim is trivial.
Assume henceforth that $\varepsilon_0>0$.
For each $z\in[a,b]$, convexity gives the chord inequality
\[
\phi(z)\le \frac{b-z}{b-a}\,\phi(a)+\frac{z-a}{b-a}\,\phi(b).
\]
Integrating and using $\int z\,\mu(dz)=1$ gives
\[
\int\phi(z)\,\mu(dz)
\le \frac{b-1}{b-a}\,\phi(a)+\frac{1-a}{b-a}\,\phi(b)
=\int \phi(z)\,\mu^\star(dz).
\qedhere
\]
\end{proof}

\begin{remark}[Martingale coupling characterization of convex order]\label{rem:strassen}
We use the following classical fact due to Strassen: if $X$ and $Y$ are integrable real random variables with the same mean, then $X\preceq_{cx}Y$ if and only if there exists a coupling $(X,Y)$ such that
\[
\E[Y\mid X]=X.
\]
See \cite{strassen}.
\end{remark}

\begin{lemma}[Convex order is preserved by averaging]\label{lem:cx-average}
Let $Z\preceq_{cx} Z^\star$ be integrable real random variables with the same mean. Then for every $n\ge 1$,
\[
\frac1n\sum_{m=1}^n Z_m\preceq_{cx}\frac1n\sum_{m=1}^n Z_m^\star,
\]
where $\{Z_m\}$ and $\{Z_m^\star\}$ are i.i.d.\ copies of $Z$ and $Z^\star$.
\end{lemma}

\begin{proof}
By Strassen's theorem for convex order, there exists a coupling $(Z,Z^\star)$ such that
\[
\E[Z^\star\mid Z]=Z.
\]
Take i.i.d.\ copies $(Z_m,Z_m^\star)$ of this coupling. Then
\[
\E\Bigl[\frac1n\sum_{m=1}^n Z_m^\star\ \Big|\ Z_1,\dots,Z_n\Bigr]=\frac1n\sum_{m=1}^n Z_m.
\]
Therefore, for every convex $\phi$,
\[
\E\phi\Bigl(\frac1n\sum_{m=1}^n Z_m\Bigr)
\le
\E\phi\Bigl(\frac1n\sum_{m=1}^n Z_m^\star\Bigr)
\]
by conditional Jensen. This is exactly the asserted convex order.
\end{proof}

\begin{theorem}[Universal privacy envelope]\label{thm:privacy-envelope}
Fix $\varepsilon_0\ge 0$ and $n\ge 1$.

\begin{enumerate}[label=\textup{(\roman*)}]
\item Let $f:[0,\infty)\to\R$ be convex with $f(1)=0$. Then
\[
\sup_{d\ge 2}\ \sup_{\rho\in\fM_d:\,W_\rho\ \varepsilon_0\text{-LDP}}\ \sup_{i\neq j}
D_f\bigl(Q_{1,n}^{ij}\,\|\,Q_{0,n}^{ij}\bigr)
=
D_f\bigl(Q_{1,n}^{\mathrm{BRR}}\,\|\,Q_{0,n}^{\mathrm{BRR}}\bigr),
\]
where the right-hand side is the one-step shuffled neighboring pair for binary randomized response with local parameter $\varepsilon_0$.

\item For every $\alpha\ge 1$,
\[
\sup_{d\ge 2}\ \sup_{\rho\in\fM_d:\,W_\rho\ \varepsilon_0\text{-LDP}}\ \sup_{i\neq j} H_\alpha\bigl(Q_{1,n}^{ij},Q_{0,n}^{ij}\bigr)
=
H_\alpha\bigl(Q_{1,n}^{\mathrm{BRR}},Q_{0,n}^{\mathrm{BRR}}\bigr),
\]
and likewise
\[
\sup_{d\ge 2}\ \sup_{\rho\in\fM_d:\,W_\rho\ \varepsilon_0\text{-LDP}}\ \sup_{i\neq j} H_\alpha\bigl(Q_{0,n}^{ij},Q_{1,n}^{ij}\bigr)
=
H_\alpha\bigl(Q_{0,n}^{\mathrm{BRR}},Q_{1,n}^{\mathrm{BRR}}\bigr).
\]
Thus binary randomized response universally maximizes the exact $(\varepsilon,\delta)$ privacy profile pointwise.
\end{enumerate}
\end{theorem}

\begin{proof}
If $\varepsilon_0=0$, then every feasible local channel has identical rows, $\mu^\star=\delta_1$, and both statements are immediate. Thus assume $\varepsilon_0>0$.

Fix a feasible triple $(d,\rho,i,j)$. Since $W_\rho$ is $\varepsilon_0$-LDP,
\[
e^{-\varepsilon_0}\le \frac{dW_\rho(\cdot\mid j)}{dW_\rho(\cdot\mid i)}\le e^{\varepsilon_0}
\qquad W_\rho(\cdot\mid i)\text{-a.s.}
\]
Hence the one-user LR law $\mu_{ij}=\Law(L_{ij})$ belongs to $\cM_{\varepsilon_0}$.
By \cref{prop:pairwise-lr}, under $Q_{0,n}^{ij}$ the shuffled likelihood ratio is
\[
\bar L_n:=\frac1n\sum_{m=1}^n L_{ij,m}.
\]
Therefore
\[
D_f(Q_{1,n}^{ij}\|Q_{0,n}^{ij})=\E f(\bar L_n).
\]
By \cref{lem:cx-max,lem:cx-average},
\[
\E f(\bar L_n)\le \E f(\bar L_n^\star),
\]
where $\bar L_n^\star$ is the average of $n$ i.i.d.\ draws from $\mu^\star$. The law $\mu^\star$ is realized by binary randomized response, so the right-hand side is exactly $D_f(Q_{1,n}^{\mathrm{BRR}}\|Q_{0,n}^{\mathrm{BRR}})$. This proves (i).

For (ii), use the convex functions
\[
f_\alpha(t):=(t-\alpha)_+,
\qquad
 g_\alpha(t):=(1-\alpha t)_+.
\]
Then
\[
H_\alpha(Q_{1,n}^{ij},Q_{0,n}^{ij})=\E f_\alpha(\bar L_n),
\qquad
H_\alpha(Q_{0,n}^{ij},Q_{1,n}^{ij})=\E g_\alpha(\bar L_n).
\]
Applying part (i) to $f_\alpha$ and $g_\alpha$ gives the result.
\end{proof}

\begin{theorem}[Privacy rigidity]\label{thm:privacy-rigidity}
Fix $\varepsilon_0\ge 0$ and $n\ge 1$. Let $\rho\in\fM_d$, let $i\neq j$, and assume that $W_\rho$ is $\varepsilon_0$-LDP. Suppose that for every $\alpha\ge 1$,
\[
H_\alpha\bigl(Q_{1,n}^{ij},Q_{0,n}^{ij}\bigr)
=
H_\alpha\bigl(Q_{1,n}^{\mathrm{BRR}},Q_{0,n}^{\mathrm{BRR}}\bigr)
\]
and
\[
H_\alpha\bigl(Q_{0,n}^{ij},Q_{1,n}^{ij}\bigr)
=
H_\alpha\bigl(Q_{0,n}^{\mathrm{BRR}},Q_{1,n}^{\mathrm{BRR}}\bigr).
\]
Then the one-user pairwise likelihood-ratio law satisfies
\[
\mu_{ij}=\mu^\star.
\]
Consequently, after merging outputs according to the two values of the pairwise likelihood ratio, the two-row reduction of $(W_\rho(\cdot\mid i),W_\rho(\cdot\mid j))$ is conditionally equivalent to binary randomized response with local parameter $\varepsilon_0$.
\end{theorem}

\begin{proof}
If $\varepsilon_0=0$, then the local constraint forces $W_\rho(\cdot\mid i)=W_\rho(\cdot\mid j)$, so $\mu_{ij}=\delta_1=\mu^\star$ and there is nothing to prove. Thus assume $\varepsilon_0>0$.

Let $\bar\mu_n$ be the law of
\[
\bar L_n:=\frac1n\sum_{m=1}^n L_{ij,m},
\]
where $L_{ij,1},\dots,L_{ij,n}$ are i.i.d.\ with law $\mu_{ij}$, and let $\bar\mu_n^\star$ be the corresponding averaged law built from $\mu^\star$. By \cref{prop:pairwise-lr},
\[
H_\alpha(Q_{1,n}^{ij},Q_{0,n}^{ij})=\int (t-\alpha)_+\,\bar\mu_n(dt),
\]
\[
H_\alpha(Q_{0,n}^{ij},Q_{1,n}^{ij})=\int (1-\alpha t)_+\,\bar\mu_n(dt),
\]
and likewise for $\bar\mu_n^\star$.

Define the stop-loss transform
\[
\mathsf S_{\bar\mu}(\alpha):=\int (t-\alpha)_+\,\bar\mu(dt),\qquad \alpha\in\R.
\]
Since both $\bar\mu_n$ and $\bar\mu_n^\star$ are compactly supported on $[e^{-\varepsilon_0},e^{\varepsilon_0}]$, their stop-loss transforms are convex and Lipschitz on $\R$. The right derivative of a convex function exists at every point, and for the stop-loss transform it satisfies
\[
\frac{d^+}{d\alpha}\mathsf S_{\bar\mu}(\alpha)=-\bar\mu((\alpha,\infty))
\qquad \forall\,\alpha\in\R.
\]
Therefore the equality
\[
\mathsf S_{\bar\mu_n}(\alpha)=\mathsf S_{\bar\mu_n^\star}(\alpha)\qquad \forall\,\alpha\ge 1
\]
implies
\[
\bar\mu_n((\alpha,\infty))=\bar\mu_n^\star((\alpha,\infty))
\qquad \forall\,\alpha\ge 1,
\]
This determines the common restriction of both measures to $(1,e^{\varepsilon_0}]$; the atom at $1$ is recovered only after combining this information with the lower stop-loss half below.

Now define the lower stop-loss transform
\[
\mathsf T_{\bar\mu}(\beta):=\int (\beta-t)_+\,\bar\mu(dt),\qquad \beta\in\R.
\]
It is also convex and Lipschitz, with right derivative
\[
\frac{d^+}{d\beta}\mathsf T_{\bar\mu}(\beta)=\bar\mu((-\infty,\beta])
\qquad \forall\,\beta\in\R.
\]
Since
\[
(1-\alpha t)_+=\alpha\Bigl(\frac1\alpha-t\Bigr)_+,
\]
the equality of the reverse hockey-stick profiles is equivalent to
\[
\mathsf T_{\bar\mu_n}(\beta)=\mathsf T_{\bar\mu_n^\star}(\beta)
\qquad\forall\,\beta\in[e^{-\varepsilon_0},1].
\]
Therefore
\[
\bar\mu_n((-\infty,\beta])=\bar\mu_n^\star((-\infty,\beta])
\qquad \forall\,\beta\in[e^{-\varepsilon_0},1],
\]
which determines both measures on $[e^{-\varepsilon_0},1]$ including any atoms. Hence
\[
\bar\mu_n=\bar\mu_n^\star.
\]

Let
\[
\phi(z):=\int e^{izt}\,\mu_{ij}(dt),
\qquad
\psi(z):=\int e^{izt}\,\mu^\star(dt)
\]
be the characteristic functions of the one-user laws. Since both laws are compactly supported, $\phi$ and $\psi$ extend to entire functions on $\mathbb C$. The characteristic function of $\bar\mu_n$ is
\[
\phi_{\bar\mu_n}(z)=\phi(z/n)^n,
\]
and similarly
\[
\phi_{\bar\mu_n^\star}(z)=\psi(z/n)^n.
\]
Because $\bar\mu_n=\bar\mu_n^\star$,
\[
\phi(z/n)^n=\psi(z/n)^n\qquad \forall\, z\in\R.
\]
Equivalently,
\[
\phi(w)^n=\psi(w)^n\qquad \forall\,w\in\R.
\]
The entire function
\[
F(w):=\phi(w)^n-\psi(w)^n
\]
vanishes on $\R$ (a set with limit points), hence vanishes identically on $\mathbb C$ by the identity theorem for entire functions.

Since $\psi(0)=1$, there exists a connected neighborhood $U$ of $0$ on which $\psi$ has no zeros. On $U$ the quotient
\[
r(w):=\frac{\phi(w)}{\psi(w)}
\]
is analytic and satisfies
\[
r(w)^n=1\qquad \forall\,w\in U.
\]
The image $r(U)$ is connected and contained in the finite set of $n$th roots of unity, so $r$ is constant on $U$. Because $r(0)=1$, we have $r\equiv 1$ on $U$. Hence $\phi=\psi$ on $U$, and therefore on all of $\mathbb C$ by the identity theorem. So
\[
\mu_{ij}=\mu^\star.
\]

Finally, since $\mu^\star$ is supported on the two points $e^{-\varepsilon_0}$ and $e^{\varepsilon_0}$, the pairwise likelihood ratio
\[
L_{ij}=\frac{dW_\rho(\cdot\mid j)}{dW_\rho(\cdot\mid i)}
\]
takes exactly these two values $W_\rho(\cdot\mid i)$-almost surely. Merge outputs according to the value of $L_{ij}$. Under row $i$, the resulting binary law has probabilities
\[
\frac{e^{\varepsilon_0}-1}{e^{\varepsilon_0}-e^{-\varepsilon_0}}
\quad\text{and}\quad
\frac{1-e^{-\varepsilon_0}}{e^{\varepsilon_0}-e^{-\varepsilon_0}}.
\]
Under row $j$, the probabilities are obtained by multiplying by the likelihood ratio:
\[
e^{-\varepsilon_0}\cdot \frac{e^{\varepsilon_0}-1}{e^{\varepsilon_0}-e^{-\varepsilon_0}}
=
\frac{1}{1+e^{\varepsilon_0}},
\qquad
e^{\varepsilon_0}\cdot \frac{1-e^{-\varepsilon_0}}{e^{\varepsilon_0}-e^{-\varepsilon_0}}
=
\frac{e^{\varepsilon_0}}{1+e^{\varepsilon_0}}.
\]
Thus the merged two-row experiment is exactly binary randomized response.
\end{proof}

\begin{remark}[Classical versus new content of the rigidity proof]
The analytic core of the uniqueness step---that compactly supported probability measures are determined by the law of their $n$-fold convolution average---is a classical consequence of the identity theorem for entire characteristic functions. The novelty of \cref{thm:privacy-rigidity} is the derivation of this rigidity condition from the saturation of both directed shuffle hockey-stick envelopes at finite $n$.
\end{remark}

\section{Estimation in the anchored gauge}\label{sec:design}

\subsection{Covariance and Fisher information}

\begin{definition}[Moment matrices]
For $\rho\in\fM_d$ define
\[
\Sigma(\rho):=\int xx^\top\,\rho(dx),
\qquad
B_i(\rho):=\int (1+\gamma_i^\top x)\,xx^\top\,\rho(dx).
\]
\end{definition}

\begin{proposition}[Fisher information at the uniform point]\label{prop:fisher}
Let $u\in\cH$ parameterize compositions by
\[
\theta(u):=\frac{\one}{d}+Hu.
\]
Then the one-user law is
\[
q_u(dx)=(1+u^\top x)\rho(dx).
\]
At $u=0$ the one-user score is $x$, so the $n$-sample Fisher information matrix is
\[
I_{\mathrm{unif}}(\rho)=n\,\Sigma(\rho).
\]
\end{proposition}

\begin{proof}
The density of $q_u$ with respect to $\rho$ is $1+u^\top x$. Hence
\[
\partial_{u_j}\log q_u(x)=\frac{x_j}{1+u^\top x},
\]
so at $u=0$ the score vector is exactly $x$. Therefore the one-user Fisher matrix at $u=0$ is
\[
\int xx^\top\,\rho(dx)=\Sigma(\rho),
\]
and the $n$-sample Fisher information is $n\Sigma(\rho)$.
\end{proof}


\begin{convention}[Singular covariance]\label{rem:canonical-convention}
When $\Sigma(\rho)$ is singular, the affine estimator based on $\Sigma(\rho)^{-1}$ is not defined. Throughout Sections~\ref{sec:design}--\ref{sec:raw}, we adopt the convention
\[
R^{\mathrm{iid}}_\theta(\rho)=R^{\mathrm{fc}}_\theta(\rho):=+\infty
\qquad\text{whenever }\Sigma(\rho)\text{ is singular}.
\]
All explicit formulas below are stated under invertibility assumptions, but the optimization theorems are understood with this convention. In particular, singular laws can never improve the canonical design objective.
\end{convention}

\begin{remark}[Mechanism knowledge]
The canonical estimator is mechanism-calibrated: it depends on $\rho$ through the matrix $\Sigma(\rho)$. This is standard in mechanism-design problems, because the analyst chooses the channel and therefore knows $\rho$ at estimation time.
\end{remark}



\subsection{Canonical unbiased estimator under fixed composition}

\begin{proposition}[Exact fixed-composition risk]\label{prop:fixed-risk}
Let $\rho\in\fM_d$ and assume $\Sigma(\rho)$ is invertible. For a deterministic composition $\theta\in\Delta_d$ with integer counts $n\theta_i$, let $X_1,\dots,X_n$ be independent outputs with exactly $n\theta_i$ draws from row $i$.
Define
\[
\widehat\theta_\rho:=\frac{\one}{d}+H\Sigma(\rho)^{-1}\bar X,
\qquad
\bar X:=\frac1n\sum_{m=1}^n X_m.
\]
Then $\widehat\theta_\rho$ is unbiased and
\[
R^{\mathrm{fc}}_\theta(\rho):=\E_\theta\|\widehat\theta_\rho-\theta\|_2^2
=\frac1n\sum_{i=1}^d \theta_i\Bigl[\tr\bigl(\Sigma^{-1}B_i\Sigma^{-1}\bigr)-\frac{d-1}{d}\Bigr].
\]
Consequently:
\begin{enumerate}[label=\textup{(\roman*)}]
\item $R^{\mathrm{fc}}_\theta(\rho)$ is affine in $\theta$;
\item $\sup_{\theta\in\Delta_d} R^{\mathrm{fc}}_\theta(\rho)=\max_i R^{\mathrm{fc}}_{e_i}(\rho)$;
\item
\[
\frac1d\sum_{i=1}^d R^{\mathrm{fc}}_{e_i}(\rho)=\frac1n\Bigl[\tr\Sigma(\rho)^{-1}-\frac{d-1}{d}\Bigr].
\]
\end{enumerate}
\end{proposition}

\begin{proof}
Under row $i$,
\[
\E_i[X]=\int x(1+\gamma_i^\top x)\,\rho(dx)=\Sigma\gamma_i.
\]
Hence under the fixed composition $\theta$,
\[
\E_\theta[\bar X]=\Sigma\sum_{i=1}^d \theta_i\gamma_i=\Sigma h_\theta.
\]
Since $Hh_\theta=\theta-\one/d$, we obtain
\[
\E_\theta[\widehat\theta_\rho]=\frac{\one}{d}+H\Sigma^{-1}\Sigma h_\theta=\theta,
\]
so the estimator is unbiased.

The row-$i$ covariance of a single draw is
\[
\Cov_i(X)=B_i-\Sigma\gamma_i\gamma_i^\top\Sigma.
\]
Since the users are independent and exactly $n\theta_i$ come from row $i$,
\[
\Cov_\theta(\bar X)=\frac1n\sum_{i=1}^d\theta_i\bigl(B_i-\Sigma\gamma_i\gamma_i^\top\Sigma\bigr).
\]
Because $H$ is an isometry from $\R^{d-1}$ to $T_d$, the squared error is
\[
\|\widehat\theta_\rho-\theta\|_2^2=\bigl\|\Sigma^{-1}(\bar X-\E_\theta\bar X)\bigr\|_2^2.
\]
Taking expectations yields
\[
\begin{aligned}
R^{\mathrm{fc}}_\theta(\rho)
&=\tr\Bigl(\Sigma^{-1}\Cov_\theta(\bar X)\Sigma^{-1}\Bigr)\\
&=\frac1n\sum_{i=1}^d\theta_i\Bigl[\tr(\Sigma^{-1}B_i\Sigma^{-1})-\tr(\gamma_i\gamma_i^\top)\Bigr].
\end{aligned}
\]
Since $\|\gamma_i\|_2^2=(d-1)/d$, the formula follows.

Affineness in $\theta$ is immediate. Since an affine function attains its maximum over the simplex at a vertex, (ii) follows. For (iii), note that
\[
\frac1d\sum_{i=1}^d B_i
=\int \frac1d\sum_{i=1}^d(1+\gamma_i^\top x)xx^\top\,\rho(dx)
=\Sigma
\]
because $\sum_i\gamma_i=0$. Averaging the vertex formula therefore gives the identity.
\end{proof}

\begin{lemma}[Exchangeable vanishing]\label{lem:exch-vanish}
Let $\rho\in\fM_d$ be exchangeable. Then for every integrable scalar function $g:\cX_d\to\R$ satisfying $g(P_\pi x)=g(x)$ for all $\pi\in S_d$,
\[
\int x\,g(x)\,\rho(dx)=0.
\]
\end{lemma}

\begin{proof}
Set $m_g:=\int x\,g(x)\,\rho(dx)$. By exchangeability of $\rho$, the vector $m_g$ satisfies $P_\pi m_g=m_g$ for every $\pi\in S_d$. Since the fixed subspace of the permutation action on $T_d$ is $\{0\}$, one has $m_g=0$.
\end{proof}

\begin{proposition}[The i.i.d.\ risk formula]\label{prop:iid-risk}
Let $\rho\in\fM_d$ with invertible $\Sigma=\Sigma(\rho)$, and suppose the user inputs are i.i.d.\ from $\theta\in\Delta_d$. Then the canonical estimator
\[
\widehat\theta_\rho=\frac{\one}{d}+H\Sigma^{-1}\bar X
\]
has risk
\[
R^{\mathrm{iid}}_\theta(\rho)
=\frac1n\int (1+h_\theta^\top x)\,x^\top\Sigma^{-2}x\,\rho(dx)-\frac{\|h_\theta\|_2^2}{n}.
\]
If $\rho$ is exchangeable, then by \cref{lem:exch-vanish} and \cref{thm:symmetrization}(iii), the cross term vanishes and therefore
\[
R^{\mathrm{iid}}_\theta(\rho)=\frac1n\tr\Sigma^{-1}-\frac{\|h_\theta\|_2^2}{n}.
\]
In particular, for exchangeable $\rho$ the maximum is attained at the uniform composition $\theta=\one/d$.
\end{proposition}

\begin{proof}
Under i.i.d.\ sampling from $\theta$, one draw has law $q_\theta(dx)=(1+h_\theta^\top x)\rho(dx)$. Its mean is
\[
\E_\theta[X]=\Sigma h_\theta,
\]
and its second moment matrix is
\[
M_\theta:=\int (1+h_\theta^\top x)\,xx^\top\,\rho(dx).
\]
Therefore
\[
\Cov_\theta(\bar X)=\frac1n\bigl(M_\theta-\Sigma h_\theta h_\theta^\top\Sigma\bigr).
\]
As in the proof of \cref{prop:fixed-risk},
\[
R^{\mathrm{iid}}_\theta(\rho)=\tr\Bigl(\Sigma^{-1}\Cov_\theta(\bar X)\Sigma^{-1}\Bigr),
\]
which is the stated formula.

If $\rho$ is exchangeable, then exchangeability implies $\Sigma(\rho)=\sigma^2 I_{d-1}$ for some $\sigma^2>0$ (by \cref{thm:symmetrization}(iii)), so $x^\top\Sigma^{-2}x=\sigma^{-4}\|x\|^2$ is permutation-invariant. Applying \cref{lem:exch-vanish} with $g(x)=x^\top\Sigma^{-2}x$ gives $\int x\,(x^\top\Sigma^{-2}x)\,\rho(dx)=0$, and hence
\[
\int (h_\theta^\top x)\,x^\top\Sigma^{-2}x\,\rho(dx)=0.
\]
Thus
\[
\int (1+h_\theta^\top x)x^\top\Sigma^{-2}x\,\rho(dx)=\int x^\top\Sigma^{-2}x\,\rho(dx)=\tr\Sigma^{-1}.
\]
The displayed formula follows, and since the second term is nonnegative, the maximum is attained at $h_\theta=0$, i.e. at $\theta=\one/d$.
\end{proof}

\begin{definition}[Canonical pairwise $\chi^2$ budget]\label{def:budget}
For $\rho\in\fM_d$, define
\[
\chi_*(\rho):=\max_{i\neq j}\int \frac{\bigl((\gamma_j-\gamma_i)^\top x\bigr)^2}{1+\gamma_i^\top x}\,\rho(dx),
\]
where the integrand is interpreted in the extended-real $\chi^2$ sense: on the set $\{1+\gamma_i^\top x=0\}$ it is taken to be $0$ when $(\gamma_j-\gamma_i)^\top x=0$ and $+\infty$ otherwise.
\end{definition}

\section{Global symmetrization}\label{sec:symmetrization}

Let $S_d$ act on $T_d$ by the restricted permutation matrices
\[
P_\pi:=H^\top \Pi H,\qquad \pi\in S_d,
\]
where $\Pi$ is the $d\times d$ permutation matrix of $\pi$.
This induces an action on $\cX_d$ because
\[
1+\gamma_i^\top(P_\pi x)=1+\gamma_{\pi^{-1}i}^\top x.
\]

\begin{definition}[Symmetrization]
For $\rho\in\fM_d$, define
\[
\bar\rho:=\frac1{d!}\sum_{\pi\in S_d}(P_\pi)_\#\rho.
\]
We call $\rho$ \emph{exchangeable} if $\bar\rho=\rho$.
\end{definition}

\begin{theorem}[Global symmetrization theorem]\label{thm:symmetrization}
For every $\rho\in\fM_d$:
\begin{enumerate}[label=\textup{(\roman*)}]
\item $\chi_*(\bar\rho)\le \chi_*(\rho)$.
\item
\[
\sup_{\theta\in\Delta_d}R^{\mathrm{fc}}_\theta(\bar\rho)
\le
\sup_{\theta\in\Delta_d}R^{\mathrm{fc}}_\theta(\rho).
\]
\item The covariance isotropizes:
\[
\Sigma(\bar\rho)=\frac{\tr\Sigma(\rho)}{d-1}\,I_{d-1}.
\]
\end{enumerate}
\end{theorem}


\begin{proof}
(i) For a fixed pair $(i,j)$, write
\[
c_{ij}(x):=\frac{((\gamma_j-\gamma_i)^\top x)^2}{1+\gamma_i^\top x}.
\]
Then
\[
c_{ij}(P_\pi x)=c_{\pi^{-1}i,\pi^{-1}j}(x).
\]
Therefore
\[
\int c_{ij}(x)\,\bar\rho(dx)=\frac1{d!}\sum_{\pi\in S_d}\int c_{\pi^{-1}i,\pi^{-1}j}(x)\,\rho(dx)\le \chi_*(\rho),
\]
and taking the maximum over $i\neq j$ gives the claim.

(ii) If $\Sigma(\rho)$ is singular, then \cref{rem:canonical-convention} gives
\[
\sup_{\theta\in\Delta_d}R^{\mathrm{fc}}_\theta(\rho)=+\infty,
\]
so there is nothing to prove. Hence assume $\Sigma(\rho)$ is invertible. By \cref{prop:fixed-risk},
\[
\sup_{\theta}R^{\mathrm{fc}}_\theta(\bar\rho)=\max_i R^{\mathrm{fc}}_{e_i}(\bar\rho).
\]
Exchangeability implies that all vertex risks are equal, so this maximum equals the average over vertices:
\[
\sup_{\theta}R^{\mathrm{fc}}_\theta(\bar\rho)=\frac1d\sum_{i=1}^d R^{\mathrm{fc}}_{e_i}(\bar\rho)=\frac1n\Bigl[\tr\Sigma(\bar\rho)^{-1}-\frac{d-1}{d}\Bigr].
\]
By part (iii),
\[
\tr\Sigma(\bar\rho)^{-1}=\frac{(d-1)^2}{\tr\Sigma(\rho)}.
\]
Now let $\lambda_1,\dots,\lambda_{d-1}$ be the eigenvalues of $\Sigma(\rho)$. The harmonic-arithmetic mean inequality gives
\[
\frac{(d-1)^2}{\tr\Sigma(\rho)}=\frac{(d-1)^2}{\sum_k \lambda_k}\le \sum_k \frac{1}{\lambda_k}=\tr\Sigma(\rho)^{-1}.
\]
Using again \cref{prop:fixed-risk}(iii),
\[
\frac1n\Bigl[\tr\Sigma(\rho)^{-1}-\frac{d-1}{d}\Bigr]=\frac1d\sum_{i=1}^d R^{\mathrm{fc}}_{e_i}(\rho)\le \sup_\theta R^{\mathrm{fc}}_\theta(\rho).
\]
Combining the last three displays proves (ii).

(iii) Set
\[
\widetilde\Sigma:=H\Sigma(\rho)H^\top,
\qquad
\widetilde\Sigma_{\mathrm{sym}}:=\frac1{d!}\sum_{\pi\in S_d}\Pi \widetilde\Sigma \Pi^\top,
\]
where $\Pi$ is the permutation matrix of $\pi$. Since $H^\top \Pi H=P_\pi$, one has
\[
\widetilde\Sigma_{\mathrm{sym}}=H\Sigma(\bar\rho)H^\top.
\]
The matrix $\widetilde\Sigma_{\mathrm{sym}}$ commutes with every permutation matrix. Therefore all of its diagonal entries are equal to some $\alpha$ and all of its off-diagonal entries are equal to some $\beta$, so
\[
\widetilde\Sigma_{\mathrm{sym}}=\alpha I_d+\beta \one\one^\top.
\]
Since
\[
\widetilde\Sigma_{\mathrm{sym}}=H\Sigma(\bar\rho)H^\top
\qquad\text{and}\qquad
H^\top\one=0,
\]
we have
\[
\widetilde\Sigma_{\mathrm{sym}}\one=H\Sigma(\bar\rho)H^\top\one=0.
\]
Hence
\[
0=\widetilde\Sigma_{\mathrm{sym}}\one=(\alpha+d\beta)\one,
\]
so $\beta=-\alpha/d$. The restriction of $\widetilde\Sigma_{\mathrm{sym}}$ to $T_d$ is therefore multiplication by $\alpha$, and
\[
\Sigma(\bar\rho)=\alpha I_{d-1}.
\]
Taking traces gives
\[
(d-1)\alpha=\tr\Sigma(\bar\rho)=\tr\Sigma(\rho),
\]
which proves the formula.
\end{proof}



\begin{corollary}[I.i.d. symmetrization]\label{cor:iid-symm}
For every $\rho\in\fM_d$,
\[
\sup_{\theta\in\Delta_d}R^{\mathrm{iid}}_\theta(\bar\rho)
\le
\sup_{\theta\in\Delta_d}R^{\mathrm{iid}}_\theta(\rho).
\]
Moreover,
\[
\sup_{\theta\in\Delta_d}R^{\mathrm{iid}}_\theta(\bar\rho)
=
\frac{(d-1)^2}{n\,\tr\Sigma(\rho)},
\]
with the convention that the right-hand side is $+\infty$ when $\tr\Sigma(\rho)=0$.
\end{corollary}

\begin{proof}
If $\tr\Sigma(\rho)=0$, then $\rho=\delta_0$, so $\bar\rho=\delta_0$ as well. By \cref{rem:canonical-convention}, both sides of the inequality are $+\infty$, and the stated formula is interpreted in the same way. Hence assume $\tr\Sigma(\rho)>0$.

By \cref{prop:iid-risk}, exchangeability of $\bar\rho$, and \cref{thm:symmetrization}(iii),
\[
\sup_{\theta\in\Delta_d}R^{\mathrm{iid}}_\theta(\bar\rho)
=
\frac1n\tr\Sigma(\bar\rho)^{-1}
=
\frac{(d-1)^2}{n\,\tr\Sigma(\rho)}.
\]
If $\Sigma(\rho)$ is singular, then \cref{rem:canonical-convention} gives
\[
\sup_{\theta\in\Delta_d}R^{\mathrm{iid}}_\theta(\rho)=+\infty,
\]
so the inequality is automatic. Otherwise, if $\lambda_1,\dots,\lambda_{d-1}$ are the eigenvalues of $\Sigma(\rho)$, the harmonic-arithmetic mean inequality yields
\[
\frac{(d-1)^2}{\tr\Sigma(\rho)}\le \tr\Sigma(\rho)^{-1}.
\]
Now \cref{prop:iid-risk} gives
\[
R^{\mathrm{iid}}_{\one/d}(\rho)=\frac1n\tr\Sigma(\rho)^{-1}\le \sup_{\theta\in\Delta_d}R^{\mathrm{iid}}_\theta(\rho),
\]
which proves the claim.
\end{proof}


\section{Exchangeable laws, trace caps, and augmented randomized response}\label{sec:tracecap}

\subsection{Template formulas for exchangeable laws}

\begin{lemma}[Signal and pairwise budgets in template coordinates]\label{lem:template-formulas}
Let $\rho\in\fM_d$ and let $\nu:=a_\#\rho\in\cP(\cK_d)$ be its template law. For $a\in\cK_d$ write
\[
A(a):=\sum_{i=1}^d \frac1{a_i}\in[ d,\infty],
\qquad
B(a):=\sum_{i=1}^d a_i^2\in[d,d^2],
\]
with the convention $A(a)=\infty$ if some $a_i=0$, and set
\[
c(a):=\frac{A(a)B(a)-d^2}{d(d-1)}.
\]
Then:
\begin{enumerate}[label=\textup{(\roman*)}]
\item
\[
\tr\Sigma(\rho)=\int_{\cK_d} (B(a)-d)\,\nu(da).
\]
\item For each ordered pair $i\neq j$,
\[
\chi_{ij}(\rho):=\int_{\cK_d}\frac{(a_j-a_i)^2}{a_i}\,\nu(da)
\]
is the directed pairwise $\chi^2$ divergence from row $i$ to row $j$ (with the same extended-real convention $0/0:=0$, $c/0:=+\infty$ for $c>0$), and
\[
\frac1{d(d-1)}\sum_{i\neq j}\chi_{ij}(\rho)=\int_{\cK_d} c(a)\,\nu(da).
\]
\item If $\rho$ is exchangeable, then every ordered pair has the same value and hence
\[
\chi_*(\rho)=\int_{\cK_d} c(a)\,\nu(da).
\]
\end{enumerate}
\end{lemma}

\begin{proof}
(i) By \cref{lem:polytope-template},
\[
\tr\Sigma(\rho)=\int \|x\|_2^2\,\rho(dx)=\int \|a-\one\|_2^2\,\nu(da)=\int (B(a)-d)\,\nu(da).
\]

(ii) The first identity is simply the definition of the directed $\chi^2$ divergence in template coordinates. Averaging over all ordered pairs,
\[
\frac1{d(d-1)}\sum_{i\neq j}\chi_{ij}(\rho)
=
\int \frac1{d(d-1)}\sum_{i\neq j}\frac{(a_j-a_i)^2}{a_i}\,\nu(da).
\]
Now
\[
\sum_{j\neq i}(a_j-a_i)^2
=\sum_{j=1}^d a_j^2-2a_i\sum_{j=1}^d a_j+d a_i^2
=B(a)-2d a_i+d a_i^2,
\]
so dividing by $a_i$ and summing in $i$ gives
\[
\sum_{i\neq j}\frac{(a_j-a_i)^2}{a_i}
=\sum_{i=1}^d\left(\frac{B(a)}{a_i}-2d+d a_i\right)
=A(a)B(a)-d^2.
\]
This proves the average formula.

(iii) If $\rho$ is exchangeable, then each ordered pair $(i,j)$ is obtained from every other ordered pair by a permutation of coordinates. Hence all directed pairwise $\chi^2$ divergences are equal. Their common value is their average, which is the integral of $c(a)$ from part (ii).
\end{proof}

\begin{lemma}[Pointwise slope inequality]\label{lem:slope}
Assume $d\ge 3$. Let $a\in\cK_d$ and define $A(a),B(a)$ as above. Then
\[
\frac{B(a)-d}{A(a)B(a)-d^2}\le \frac{1}{d+2\sqrt{d-1}},
\]
with the conventions that the left-hand side is $0$ when $a=\one$ (both numerator and denominator vanish) and $0$ whenever some $a_i=0$ (since then $A(a)=\infty$ while $B(a)-d>0$). Equivalently, in cross-multiplied form:
\[
\bigl(B(a)-d\bigr)\bigl(d+2\sqrt{d-1}\bigr)\le A(a)B(a)-d^2
\qquad\text{for all }a\in\cK_d,
\]
where the right-hand side is interpreted in $[0,\infty]$; if some $a_i=0$, then $A(a)=\infty$ and the inequality holds strictly in the extended-real sense. The inequality holds with equality at $a=\one$ (trivially, both sides are zero) and is strict for all $a\notin\{\one\}\cup\{\text{permutations of }a_\star\}$. The nontrivial equality holds if and only if $a$ is a permutation of the singleton GRR template
\[
a_\star:=\left(\frac{d\sqrt{d-1}}{\sqrt{d-1}+d-1},\frac{d}{\sqrt{d-1}+d-1},\dots,\frac{d}{\sqrt{d-1}+d-1}\right).
\]
\end{lemma}


\begin{proof}
If some $a_i=0$, then $A(a)=\infty$ and, since $\sum_i a_i=d$, one has $B(a)>d$; thus the stated convention indeed gives the left-hand side as $0$, and the inequality is trivial.
If $a=\one$, then $B(a)=d$ and both sides of the cross-multiplied inequality are $0$, so the claim is immediate.
Hence it remains to treat the case $a_i>0$ for all $i$ and $a\neq\one$; equivalently, $B\in(d,d^2)$.
Set
\[
\cF_B:=\Bigl\{a\in(0,\infty)^d:\sum_{i=1}^d a_i=d,\ \sum_{i=1}^d a_i^2=B\Bigr\},
\qquad B\in[d,d^2).
\]
For fixed $B$, maximizing
\[
\frac{B-d}{A(a)B-d^2}
\]
is equivalent to minimizing $A(a)=\sum_i a_i^{-1}$ over $\cF_B$.
Because $A(a)\to \infty$ when any coordinate tends to $0$, the minimum is attained in the interior of $\cF_B$. Introduce the Lagrangian
\[
L(a;\mu_0,\nu_0)=\sum_{i=1}^d a_i^{-1}+\mu_0\Bigl(\sum_i a_i-d\Bigr)+\nu_0\Bigl(\sum_i a_i^2-B\Bigr).
\]
At a minimizer,
\[
-\frac{1}{a_i^2}+\mu_0+2\nu_0 a_i=0
\qquad (i=1,\dots,d).
\]
Thus each coordinate is a positive root of
\[
g(x):=-x^{-2}+\mu_0+2\nu_0 x.
\]
If $\nu_0\ge 0$, then $g'(x)=2x^{-3}+2\nu_0>0$, so $g$ has at most one positive root. If $\nu_0<0$, then $g'$ is strictly decreasing and can vanish at most once, so $g$ has at most two positive zeros. Hence every minimizer has at most two distinct positive coordinates.

Therefore any extremizer can be written as
\[
a=(\underbrace{\alpha,\dots,\alpha}_{s\text{ times}},\underbrace{\beta,\dots,\beta}_{d-s\text{ times}})
\]
with $1\le s\le d-1$ and $\alpha>\beta>0$. Write $\lambda:=\alpha/\beta>1$. From the sum constraint,
\[
\alpha=\frac{d\lambda}{d+s(\lambda-1)},
\qquad
\beta=\frac{d}{d+s(\lambda-1)}.
\]
A direct substitution gives
\[
B-d=\frac{d\,s(d-s)(\lambda-1)^2}{(d+s(\lambda-1))^2}
\]
and
\[
AB-d^2=\frac{d\,s(d-s)(\lambda-1)^2(\lambda+1)}{\lambda\,(d+s(\lambda-1))}.
\]
Therefore the slope reduces to
\[
r_s(\lambda):=\frac{B-d}{AB-d^2}=\frac{\lambda}{(\lambda+1)(d+s(\lambda-1))}.
\]
Differentiating,
\[
r_s'(\lambda)=\frac{d-s-s\lambda^2}{(\lambda+1)^2(d+s(\lambda-1))^2}.
\]
If $1\le s<d/2$, the unique critical point is
\[
\lambda_s^\star=\sqrt{\frac{d-s}{s}}>1,
\]
and at this point
\[
\max_{\lambda>1} r_s(\lambda)=\frac{1}{d+2\sqrt{s(d-s)}}.
\]
If $s=d/2$ (necessarily with $d$ even), then $r_s'(\lambda)\le 0$ for every $\lambda\ge 1$ and
\[
\sup_{\lambda>1} r_s(\lambda)=\lim_{\lambda\downarrow 1} r_s(\lambda)=\frac{1}{2d}
<
\frac{1}{d+2\sqrt{d-1}}
\qquad (d\ge 3).
\]
If $s>d/2$, then again $d-s-s\lambda^2<0$ for every $\lambda\ge 1$, so $r_s$ is strictly decreasing on $[1,\infty)$ and
\[
\sup_{\lambda>1} r_s(\lambda)=\lim_{\lambda\downarrow 1} r_s(\lambda)=\frac{1}{2d}
<
\frac{1}{d+2\sqrt{d-1}}.
\]
Since $\sqrt{s(d-s)}$ is minimized over $1\le s<d/2$ exactly at $s=1$, the overall maximum is
\[
\frac{1}{d+2\sqrt{d-1}},
\]
attained exactly when $s=1$ and
\[
\lambda=\sqrt{d-1}.
\]
Equivalently, one coordinate takes the high value and the remaining $d-1$ coordinates take the low value, up to permutation of coordinate labels. The corresponding template is precisely $a_\star$ up to permutation.
\end{proof}


\subsection{Trace cap and explicit optimizer}

\begin{definition}[Augmented randomized response]\label{def:aug-grr}
For $\lambda>1$ and $p\in[0,1]$, define
\[
t_\lambda:=\frac{d(\lambda-1)}{d+\lambda-1}
\]
and the anchored law
\[
\rho_{p,\lambda}:=(1-p)\,\delta_0+\frac{p}{d}\sum_{i=1}^d \delta_{t_\lambda\gamma_i}.
\]
The corresponding channel $W_{\rho_{p,\lambda}}$ is the \emph{augmented randomized response} mechanism with activation probability $p$ and aggressive parameter $\lambda$.
\end{definition}

\begin{lemma}[Explicit formulas for augmented randomized response]\label{lem:aug-formulas}
For $\rho_{p,\lambda}$ as above,
\[
\Sigma(\rho_{p,\lambda})=\frac{p d(\lambda-1)^2}{(d+\lambda-1)^2}I_{d-1},
\]
\[
\chi_*(\rho_{p,\lambda})=p\,C_1(\lambda),
\qquad
C_1(\lambda):=\frac{(\lambda-1)^2(\lambda+1)}{\lambda(d+\lambda-1)}.
\]
In particular, for
\[
\lambda_*=\sqrt{d-1},
\qquad
C_*(d):=C_1(\lambda_*),
\]
one has
\[
\tr\Sigma(\rho_{p,\lambda_*})=\frac{d(d-1)}{d+2\sqrt{d-1}}\,p\,C_*(d).
\]
\end{lemma}

\begin{proof}
Since $\sum_i \gamma_i\gamma_i^\top=I_{d-1}$,
\[
\Sigma(\rho_{p,\lambda})=\frac{p}{d}\sum_{i=1}^d t_\lambda^2\gamma_i\gamma_i^\top=\frac{p t_\lambda^2}{d}I_{d-1},
\]
which is the first formula.

The template corresponding to $t_\lambda\gamma_i$ has one coordinate
\[
\alpha=1+t_\lambda\Bigl(1-\frac1d\Bigr)=\frac{d\lambda}{d+\lambda-1}
\]
and $d-1$ coordinates
\[
\beta=1-\frac{t_\lambda}{d}=\frac{d}{d+\lambda-1}.
\]
This is exactly the singleton two-level template with ratio $\lambda$, so the budget formula is the $s=1$ case of the computation in the proof of \cref{lem:slope}:
\[
\chi_*(\rho_{p,\lambda})=p\,\frac{(\lambda-1)^2(\lambda+1)}{\lambda(d+\lambda-1)}.
\]
Finally,
\[
\frac{\tr\Sigma(\rho_{p,\lambda})}{\chi_*(\rho_{p,\lambda})}=\frac{d(d-1)\lambda}{(\lambda+1)(d+\lambda-1)}.
\]
At $\lambda=\sqrt{d-1}$ this ratio equals $d(d-1)/(d+2\sqrt{d-1})$.
\end{proof}

\begin{theorem}[Exact trace cap in the low-budget regime]\label{thm:trace-cap}
Assume $d\ge 3$. Then every anchored law $\rho\in\fM_d$ satisfies
\[
\tr\Sigma(\rho)\le \frac{d(d-1)}{d+2\sqrt{d-1}}\,\chi_*(\rho).
\]
Consequently,
\[
\sup_{\rho:\,\chi_*(\rho)\le C}\tr\Sigma(\rho)\le \frac{d(d-1)}{d+2\sqrt{d-1}}\,C
\qquad \text{for all }C\ge 0.
\]
If $0\le C\le C_*(d)$, equality is attained by augmented randomized response with
\[
\lambda_*=\sqrt{d-1},\qquad p=C/C_*(d),
\]
and the optimizing anchored law is exactly $\rho_{p,\lambda_*}$; equivalently, the corresponding channel is unique up to conditionally identical refinements.
\end{theorem}

\begin{proof}
Let $\nu=a_\#\rho$ be the template law of $\rho$. By \cref{lem:template-formulas,lem:slope},
\[
\tr\Sigma(\rho)=\int (B(a)-d)\,\nu(da)
\le \frac{d(d-1)}{d+2\sqrt{d-1}}\int c(a)\,\nu(da).
\]
Using \cref{lem:template-formulas}(ii),
\[
\int c(a)\,\nu(da)
=
\frac1{d(d-1)}\sum_{i\neq j}\chi_{ij}(\rho)
\le
\chi_*(\rho).
\]
Combining the last two displays proves
\[
\tr\Sigma(\rho)\le \frac{d(d-1)}{d+2\sqrt{d-1}}\,\chi_*(\rho).
\]

For attainability in the low-budget regime, take $\rho_{p,\lambda_*}$ with $p=C/C_*(d)$. By \cref{lem:aug-formulas},
\[
\chi_*(\rho_{p,\lambda_*})=C,
\qquad
\tr\Sigma(\rho_{p,\lambda_*})=\frac{d(d-1)}{d+2\sqrt{d-1}}\,C.
\]
So equality is attained.

Now suppose $\rho$ attains equality at some $C\le C_*(d)$. Then equality must hold simultaneously in both inequalities above.

Equality in the pointwise slope inequality forces $\nu$-almost every template $a\neq\one$ to attain equality in \cref{lem:slope}. Therefore the support of $\nu$ is contained in the union of the null template $\one$ and the $d$ singleton templates
\[
a_\star^{(k)},\qquad k=1,\dots,d,
\]
obtained by placing the high coordinate of $a_\star$ in position $k$. Write
\[
\nu=(1-p)\,\delta_{\one}+\sum_{k=1}^d p_k\,\delta_{a_\star^{(k)}},
\qquad
p_k\ge 0,\quad \sum_{k=1}^d p_k=p.
\]

Let $\alpha_\star$ and $\beta_\star$ denote the high and low coordinates of $a_\star$, and set
\[
u:=\frac{(\beta_\star-\alpha_\star)^2}{\alpha_\star},
\qquad
v:=\frac{(\alpha_\star-\beta_\star)^2}{\beta_\star}.
\]
For an ordered pair $i\neq j$, only the singleton templates with high coordinate in $i$ or in $j$ contribute to $\chi_{ij}(\rho)$, hence
\[
\chi_{ij}(\rho)=u\,p_i+v\,p_j.
\]
Equality in the averaging step
\[
\frac1{d(d-1)}\sum_{i\neq j}\chi_{ij}(\rho)\le \chi_*(\rho)
\]
implies that all ordered pair budgets are equal to the common maximum $\chi_*(\rho)$; otherwise their average would be strictly smaller than their maximum. Therefore
\[
u\,p_i+v\,p_j=\chi_*(\rho)\qquad \forall\, i\neq j.
\]
Fixing $j$ and comparing two different indices $i$ and $i'$ gives $u p_i=u p_{i'}$, so all $p_i$ are equal. Hence
\[
p_i=\frac{p}{d}\qquad (i=1,\dots,d),
\]
and therefore
\[
\nu=(1-p)\,\delta_{\one}+\frac{p}{d}\sum_{k=1}^d \delta_{a_\star^{(k)}}.
\]
This is exactly the template law of augmented randomized response. Passing back through the affine bijection of \cref{lem:polytope-template} gives
\[
\rho=\rho_{p,\lambda_*}.
\]
By \cref{thm:representation}(ii)--(iii), the corresponding channel is therefore unique up to conditionally identical refinements.
\end{proof}

\subsection{Exact finite-\texorpdfstring{$n$}{n} canonical optimality under the \texorpdfstring{$\chi_*$}{chi*}-budget}

\begin{theorem}[Exact finite-$n$ canonical optimality under the $\chi_*$-budget]\label{thm:finite-canonical-chi}
Assume $d\ge 3$ and $0<C\le C_*(d)$. Then
\[
\inf_{\rho:\,\chi_*(\rho)\le C}\ \sup_{\theta\in\Delta_d}R^{\mathrm{iid}}_\theta(\rho)
=
\frac{(d-1)(d+2\sqrt{d-1})}{n d\,C},
\]
and
\[
\inf_{\rho:\,\chi_*(\rho)\le C}\ \sup_{\theta\in\Delta_d}R^{\mathrm{fc}}_\theta(\rho)
=
\frac1n\left[
\frac{(d-1)(d+2\sqrt{d-1})}{d\,C}
-\frac{d-1}{d}
\right].
\]
Both values are attained by augmented randomized response with
\[
\lambda_*=\sqrt{d-1},
\qquad
p=C/C_*(d).
\]
Moreover, the optimizing anchored law is exactly $\rho_{p,\lambda_*}$; equivalently, the corresponding channel is unique up to conditionally identical refinements.
\end{theorem}

\begin{proof}
For the i.i.d.\ risk, \cref{cor:iid-symm} shows that symmetrization cannot increase the objective. Hence
\[
\inf_{\rho:\,\chi_*(\rho)\le C}\ \sup_{\theta}R^{\mathrm{iid}}_\theta(\rho)
=
\inf_{\rho=\bar\rho:\,\chi_*(\rho)\le C}\ \sup_{\theta}R^{\mathrm{iid}}_\theta(\rho).
\]
If $\rho$ is exchangeable, then by \cref{prop:iid-risk} and \cref{thm:symmetrization}(iii),
\[
\sup_{\theta\in\Delta_d}R^{\mathrm{iid}}_\theta(\rho)
=
\frac1n\tr\Sigma(\rho)^{-1}
=
\frac{(d-1)^2}{n\,\tr\Sigma(\rho)}.
\]
Applying \cref{thm:trace-cap} gives
\[
\sup_{\theta}R^{\mathrm{iid}}_\theta(\rho)\ge
\frac{(d-1)^2}{n\,\frac{d(d-1)}{d+2\sqrt{d-1}}\,C}
=
\frac{(d-1)(d+2\sqrt{d-1})}{n d\,C}.
\]
For augmented randomized response with $p=C/C_*(d)$ and $\lambda=\lambda_*$, \cref{lem:aug-formulas} gives
\[
\Sigma=\frac{dC}{d+2\sqrt{d-1}}\,I_{d-1},
\]
hence equality holds. If a law $\rho$ (not necessarily exchangeable) attains the infimum, then $\bar\rho$ also attains it, since symmetrization does not increase the risk. Hence $\bar\rho$ is an exchangeable optimizer, which forces $\bar\rho=\rho_{p,\lambda_*}$ by the exchangeable uniqueness just established. Now $\tr\Sigma(\rho)=\tr\Sigma(\bar\rho)=K_d\,C$, where $K_d:=d(d-1)/(d+2\sqrt{d-1})$, and $\chi_*(\rho)\le C$. The trace-cap inequality (\cref{thm:trace-cap}) gives $K_d\,C=\tr\Sigma(\rho)\le K_d\,\chi_*(\rho)$, so $\chi_*(\rho)=C$ and the trace cap is saturated. By the uniqueness clause of \cref{thm:trace-cap}, $\rho=\rho_{p,\lambda_*}$.

For the fixed-composition risk, \cref{thm:symmetrization} already yields the corresponding symmetrization inequality. If $\rho$ is exchangeable, then \cref{prop:fixed-risk} and \cref{thm:symmetrization}(iii) give
\[
\sup_{\theta\in\Delta_d}R^{\mathrm{fc}}_\theta(\rho)
=
\frac1n\left[\frac{(d-1)^2}{\tr\Sigma(\rho)}-\frac{d-1}{d}\right].
\]
Applying the same trace-cap bound yields
\[
\sup_{\theta}R^{\mathrm{fc}}_\theta(\rho)\ge
\frac1n\left[
\frac{(d-1)(d+2\sqrt{d-1})}{d\,C}
-\frac{d-1}{d}
\right],
\]
and augmented randomized response attains equality by the same explicit covariance formula. Uniqueness of the optimizer follows by the same symmetrization-then-trace-cap argument as in the i.i.d.\ case.
\end{proof}

\subsection{Two-orbit reduction of the global \texorpdfstring{$\chi_*$}{chi*} frontier}

\begin{definition}[Orbit laws and the frontier curve]\label{def:orbit-frontier}
Let
\[
\Omega_d:=\{a\in\cK_d:\ a_1\ge a_2\ge \cdots \ge a_d,\ a_d>0\}.
\]
For $a\in\Omega_d$, let $\rho_a$ denote the exchangeable anchored law obtained by uniformly distributing mass on the permutation orbit of $x(a)=H^\top(a-\one)$. Define
\[
S(a):=B(a)-d,
\qquad
C(a):=\frac{A(a)B(a)-d^2}{d(d-1)},
\]
and let
\[
\Gamma_d:=\{(C(a),S(a)): a\in\Omega_d\}\subset [0,\infty)\times [0,d(d-1)].
\]
\end{definition}

\begin{lemma}[Planar Carath\'eodory reduction]\label{lem:planar-car}
Every point in the convex hull of a subset of $\R^2$ is a convex combination of at most three points of that subset.
\end{lemma}

\begin{proof}
Let $z=\sum_{r=1}^m \lambda_r z_r$ be a convex combination of points $z_r\in\R^2$ with $\lambda_r>0$ and $\sum_r\lambda_r=1$. If $m\le 3$, there is nothing to prove. Assume $m\ge 4$. Then the vectors $(1,z_r)\in\R^3$ are linearly dependent, so there exist real numbers $c_r$, not all zero, such that
\[
\sum_{r=1}^m c_r=0,
\qquad
\sum_{r=1}^m c_r z_r=0.
\]
Choose
\[
\varepsilon:=\min_{c_r>0}\frac{\lambda_r}{c_r}>0.
\]
Then
\[
\lambda_r':=\lambda_r-\varepsilon c_r\ge 0
\]
for every $r$, and equality holds for at least one index. Moreover,
\[
\sum_{r=1}^m \lambda_r'=1,
\qquad
\sum_{r=1}^m \lambda_r' z_r=z.
\]
Thus $z$ is represented as a convex combination of strictly fewer than $m$ points. Iterating reduces the number of points to at most three.
\end{proof}

\begin{theorem}[Two-orbit reduction of the global $\chi_*$ frontier]\label{thm:two-orbit}
For $C_0\ge 0$, define
\[
F_d(C_0):=\sup_{\rho:\,\chi_*(\rho)\le C_0}\tr\Sigma(\rho).
\]
Then
\[
F_d(C_0)=\sup\{s:\ (c,s)\in \conv(\Gamma_d),\ c\le C_0\}.
\]
Moreover, for every $C_0\ge 0$ there exist $a,b\in\Omega_d$ and $t\in[0,1]$ such that
\[
F_d(C_0)=t\,S(a)+(1-t)\,S(b),
\qquad
t\,C(a)+(1-t)\,C(b)= C_0.
\]
Equivalently, some optimizer of the original frontier problem is a mixture of at most two exchangeable orbit laws, and every such optimizer saturates the budget exactly.
\end{theorem}


\begin{proof}
If $C_0=0$, then $\chi_*(\rho)=0$. When $d\ge 3$, \cref{thm:trace-cap} gives $\tr\Sigma(\rho)=0$. When $d=2$, the nonnegative pairwise $\chi^2$ integrand vanishes $\rho$-a.s., which forces $(\gamma_2-\gamma_1)^\top x=0$ $\rho$-a.s.; since $\gamma_2-\gamma_1$ spans $\R^1$, this gives $\rho=\delta_0$ and $\tr\Sigma(\rho)=0$. In either case $F_d(0)=0$, realized by $\rho_{\one}$, and the two-orbit conclusion is immediate. Hence assume $C_0>0$.

Let $\rho$ be any feasible law. By \cref{thm:symmetrization}, the symmetrized law $\bar\rho$ is feasible and has the same trace, so
\[
F_d(C_0)=\sup_{\rho=\bar\rho:\,\chi_*(\rho)\le C_0}\tr\Sigma(\rho).
\]
For an exchangeable law $\rho$, let $\nu$ be its probability law on $\Omega_d$ of sorted templates. By \cref{lem:template-formulas},
\[
\chi_*(\rho)=\int_{\Omega_d} C(a)\,\nu(da),
\qquad
\tr\Sigma(\rho)=\int_{\Omega_d} S(a)\,\nu(da).
\]
Therefore the achievable exchangeable pairs are precisely the barycenters of probability measures on $\Gamma_d$, and hence they form the convex hull $\conv(\Gamma_d)$. This proves the formula for $F_d(C_0)$.

To prove existence of an optimizer, consider the equivalent optimization problem
\[
F_d(C_0)=\sup\Bigl\{\int_{\Omega_d} S(a)\,\nu(da):\ \nu\in\cP(\Omega_d),\ \int_{\Omega_d} C(a)\,\nu(da)\le C_0\Bigr\}.
\]
Let $(\nu_m)$ be a maximizing sequence. Since $B(a)\ge d$,
\[
C(a)=\frac{A(a)B(a)-d^2}{d(d-1)}\ge \frac{dA(a)-d^2}{d(d-1)}=\frac{A(a)-d}{d-1}.
\]
Hence, for every $M>0$, the level set
\[
\Omega_d(M):=\{a\in\Omega_d:\ C(a)\le M\}
\]
is compact: indeed $C(a)\le M$ implies $A(a)\le d+(d-1)M$, and since $A(a)\ge 1/a_d$, this gives
\[
a_d\ge \frac{1}{d+(d-1)M}>0.
\]
Moreover,
\[
\nu_m(\Omega_d\setminus \Omega_d(M))\le \frac{1}{M}\int_{\Omega_d} C(a)\,\nu_m(da)\le \frac{C_0}{M},
\]
so $(\nu_m)$ is tight. Because $\Omega_d$ is Polish, Prokhorov's theorem yields a subsequence, still denoted $(\nu_m)$, converging weakly to some $\nu_\star\in\cP(\Omega_d)$. Since $S$ is bounded and continuous on $\Omega_d$,
\[
\int_{\Omega_d} S(a)\,\nu_m(da)\longrightarrow \int_{\Omega_d} S(a)\,\nu_\star(da).
\]
Since $C$ is nonnegative and continuous, Portmanteau's theorem gives
\[
\int_{\Omega_d} C(a)\,\nu_\star(da)\le \liminf_{m\to\infty}\int_{\Omega_d} C(a)\,\nu_m(da)\le C_0.
\]
Thus $\nu_\star$ is feasible and optimal.

The map $F_d:[0,\infty)\to[0,d(d-1)]$ is finite, nondecreasing, and concave. Since $C_0>0$, choose a supergradient $\lambda\ge 0$ at $C_0$, so that
\[
F_d(C)\le F_d(C_0)+\lambda(C-C_0)
\qquad \forall\, C\ge 0.
\]
If $\lambda=0$, then for every $a\in\Omega_d$,
\[
S(a)\le F_d(C(a))\le F_d(C_0).
\]
Taking the supremum over $a$ gives $d(d-1)=\sup_{a\in\Omega_d}S(a)\le F_d(C_0)\le d(d-1)$, hence $F_d(C_0)=d(d-1)$. But $S(a)<d(d-1)$ for every $a\in\Omega_d$ because $a_d>0$, so no feasible probability measure can achieve expectation $d(d-1)$, a contradiction. Therefore $\lambda>0$. In particular, no optimizer can use strictly less than the full budget: if some feasible $\nu$ satisfied
\[
\int_{\Omega_d} S(a)\,\nu(da)=F_d(C_0)
\qquad\text{and}\qquad
\int_{\Omega_d} C(a)\,\nu(da)<C_0,
\]
then $F_d$ would be constant on a left-neighborhood of $C_0$, which would force $0$ to be a supergradient at $C_0$, contradicting the previous paragraph.

Define
\[
g_\lambda(a):=S(a)-\lambda C(a),\qquad a\in\Omega_d.
\]
Since $S$ is bounded and $\lambda>0$ while $C(a)\to\infty$ as $a_d\downarrow 0$, the function $g_\lambda$ attains its maximum on $\Omega_d$; let
\[
M_\lambda:=\argmax_{a\in\Omega_d} g_\lambda(a).
\]
For every $a\in\Omega_d$,
\[
S(a)\le F_d(C(a))\le F_d(C_0)+\lambda(C(a)-C_0),
\]
so
\[
g_\lambda(a)\le F_d(C_0)-\lambda C_0.
\]
On the other hand, optimality of $\nu_\star$ gives
\[
\int_{\Omega_d} g_\lambda(a)\,\nu_\star(da)
=F_d(C_0)-\lambda\int_{\Omega_d} C(a)\,\nu_\star(da)
\ge F_d(C_0)-\lambda C_0.
\]
Hence equality holds throughout. In particular,
\[
\int_{\Omega_d} g_\lambda(a)\,\nu_\star(da)=\sup_{a\in\Omega_d} g_\lambda(a)=F_d(C_0)-\lambda C_0,
\]
so $\nu_\star$ is supported on $M_\lambda$, and since $\lambda>0$ we also have
\[
\int_{\Omega_d} C(a)\,\nu_\star(da)=C_0.
\]
Thus the optimal point
\[
(C_0,F_d(C_0))
\]
lies in the convex hull of the subset
\[
\Gamma_d\cap \{(c,s): s-\lambda c = F_d(C_0)-\lambda C_0\}.
\]
By \cref{lem:planar-car}, it is a convex combination of at most three points of this subset. But all such points are collinear, since they lie on the same supporting line
\[
s-\lambda c = F_d(C_0)-\lambda C_0.
\]
Any point in the convex hull of collinear points lies on a line segment joining two of them. Therefore there exist $a,b\in\Omega_d$ and $t\in[0,1]$ such that
\[
(C_0,F_d(C_0))=t\,(C(a),S(a))+(1-t)\,(C(b),S(b)).
\]
Equivalently,
\[
F_d(C_0)=t\,S(a)+(1-t)\,S(b),
\qquad
t\,C(a)+(1-t)\,C(b)=C_0.
\]
Hence the two-orbit law
\[
\rho=t\,\rho_a+(1-t)\,\rho_b
\]
is feasible and optimal.
\end{proof}


\begin{corollary}[Concave-envelope interpretation]\label{cor:concave-envelope}
The map $C\mapsto F_d(C)$ is the upper concave envelope of the parametric curve $\Gamma_d$.
\end{corollary}

\begin{proof}
The preceding theorem identifies the achievable region with the convex hull of $\Gamma_d$ and the frontier with its upper boundary. The upper boundary of a planar convex set is exactly the graph of a concave function, namely the upper concave envelope of the generating curve.
\end{proof}

\section{Universal low-budget optimality}\label{sec:asymptotic}

\subsection{Asymptotic minimax theorem}

We now consider the i.i.d.\ model: user inputs are i.i.d.\ from an unknown composition $\theta\in\Delta_d$, each input is privatized through a channel represented by $\rho\in\fM_d$, and the observer receives the shuffled outputs. Since the privatized messages are i.i.d.\ under this input model, the shuffled ordered sample has law $q_\theta^{\otimes n}$. Equivalently, the unordered sample (the empirical measure; for finite-output channels, the histogram) is a sufficient statistic for this i.i.d.\ product family, so minimax and Bayes analyses may be carried out in the product model $\{q_\theta^{\otimes n}:\theta\in\Delta_d\}$.

\begin{theorem}[Universal low-budget optimality among all estimators]\label{thm:main-design}
Assume $d\ge 3$, and let $C_n\downarrow 0$ with $nC_n\to\infty$. Then
\[
\inf_{\rho:\,\chi_*(\rho)\le C_n}\ \inf_{\widehat\theta}\ \sup_{\theta\in\Delta_d}
\E_{\theta,\rho}\|\widehat\theta-\theta\|_2^2
=
\frac{d-1}{nd}\,\frac{d+2\sqrt{d-1}}{C_n}\,(1+o(1)).
\]
Here the outer infimum is over mechanism design (choice of $\rho$), the inner infimum is over \emph{all} estimators $\widehat\theta$ (not only the canonical one), and the supremum is over all compositions.
The upper bound is achieved, for all sufficiently large $n$, by augmented randomized response with
\[
\lambda_*=\sqrt{d-1},\qquad p_n=C_n/C_*(d),
\]
and the canonical estimator.
\end{theorem}

\begin{proof}
\textbf{Upper bound.}
Let $\rho_n=\rho_{p_n,\lambda_*}$. By \cref{lem:aug-formulas},
\[
\Sigma(\rho_n)=s_n I_{d-1},
\qquad
s_n=\frac{dC_n}{d+2\sqrt{d-1}}.
\]
Using \cref{prop:iid-risk},
\[
R^{\mathrm{iid}}_\theta(\rho_n)=\frac{d-1}{n s_n}-\frac{\|h_\theta\|_2^2}{n}.
\]
Hence the worst case is at the uniform composition and
\[
\sup_{\theta\in\Delta_d}R^{\mathrm{iid}}_\theta(\rho_n)=\frac{d-1}{n s_n}=\frac{d-1}{nd}\,\frac{d+2\sqrt{d-1}}{C_n}.
\]
This proves the matching upper bound.

\textbf{Lower bound.}
Let $m:=d-1$. Because $0\in\operatorname{int}\cH$, there exists $\delta_d>0$ such that the cube $[-\delta_d,\delta_d]^m$ is contained in $\cH$. Choose
\[
r_n:=(nC_n)^{-1/4},
\]
so that $r_n\to 0$ and, for all large $n$, $[-r_n,r_n]^m\subset\cH$.
Parameterize locally by
\[
\theta(u):=\frac{\one}{d}+Hu,
\qquad u\in[-r_n,r_n]^m.
\]
Since $H$ is an isometry from $\R^{d-1}$ to $T_d$, it suffices to lower-bound the risk for the local coordinate $u$. Given an arbitrary estimator $\widehat\theta$, define $\widehat u:=H^\top\!\left(\widehat\theta-\frac{\one}{d}\right)$. Because $H^\top H=I_{d-1}$ and $HH^\top\preceq I_d$ (since $HH^\top$ is the orthogonal projection onto $T_d$),
\[
\|\widehat\theta-\theta(u)\|_2^2\ge \|H^\top(\widehat\theta-\theta(u))\|_2^2=\|\widehat u-u\|_2^2.
\]
Therefore any Bayes lower bound for estimating $u$ applies a fortiori to the original estimator $\widehat\theta$.

Fix a one-dimensional $C^1$ density $\varphi$ on $[-1,1]$ such that $\varphi(\pm 1)=0$ and
\[
J_0:=\int_{-1}^1 \frac{(\varphi'(t))^2}{\varphi(t)}\,dt<\infty.
\]
For concreteness one may take
\[
\varphi(t)=\frac{15}{16}(1-t^2)^2\,\mathbf 1_{[-1,1]}(t).
\]
Set
\[
\varphi_{r_n}(s):=r_n^{-1}\varphi(s/r_n),
\qquad
\pi_n(du):=\prod_{j=1}^m \varphi_{r_n}(u_j)\,du_j.
\]
Then
\[
J(\pi_n):=\sum_{j=1}^m \int \frac{(\partial_j \pi_n(u))^2}{\pi_n(u)}\,du=\frac{mJ_0}{r_n^2}=o(nC_n).
\]

Now fix any feasible $\rho$ with $\chi_*(\rho)\le C_n$. For the product model $q_u^{\otimes n}$, the one-user score in direction $u_j$ is
\[
\partial_{u_j}\log q_u(x)=\frac{x_j}{1+u^\top x}.
\]
Hence
\[
I_u(\rho)=n\int \frac{xx^\top}{1+u^\top x}\,\rho(dx).
\]
By \cref{lem:polytope-template}, every $x\in\cX_d$ satisfies $\|x\|\le R_d:=\sqrt{d(d-1)}$. Since $\|u\|\le \sqrt m\,r_n$ on the support of $\pi_n$,
\[
1+u^\top x\ge 1-R_d\sqrt m\,r_n=1-o(1),
\]
uniformly in $u$ and $x$. Therefore
\[
\tr I_u(\rho)\le \frac{n}{1-R_d\sqrt m\,r_n}\,\tr\Sigma(\rho)=n(1+o(1))\tr\Sigma(\rho).
\]
By \cref{thm:trace-cap},
\[
\tr\Sigma(\rho)\le \frac{d(d-1)}{d+2\sqrt{d-1}}\,C_n=:\alpha_d C_n.
\]
Thus
\[
\int \tr I_u(\rho)\,\pi_n(du)\le n\alpha_d C_n(1+o(1)).
\]

Apply the multivariate van Trees bound from \cref{lem:vantrees} in the appendix. For every estimator $\widehat u$,
\[
\int \E_u\|\widehat u-u\|_2^2\,\pi_n(du)
\ge
\frac{m^2}{\int \tr I_u(\rho)\,\pi_n(du)+J(\pi_n)}.
\]
Combining the previous two displays gives
\[
\int \E_u\|\widehat u-u\|_2^2\,\pi_n(du)
\ge
\frac{m^2}{n\alpha_d C_n(1+o(1))}
=\frac{d-1}{nd}\,\frac{d+2\sqrt{d-1}}{C_n}(1-o(1)).
\]
Because Bayes risk is bounded above by minimax risk and $\|\widehat\theta-\theta\|_2^2\ge\|\widehat u-u\|_2^2$, we conclude that every estimator under every feasible $\rho$ obeys the same lower bound. Together with the upper bound this proves the theorem.
\end{proof}

\section{Exact design under the raw local \texorpdfstring{$\varepsilon_0$}{epsilon0} cap}\label{sec:raw}

The optimizer class for this problem---subset selection---was identified by Wang, Blocki, Li, and Jha \cite{wang-ldp} and shown to be asymptotically order-optimal by Ye and Barg \cite{ye-barg}. Pan \cite{pan-strict} recently proved strict optimality in the pure local model. The present section does not claim a new identification of the optimizer class. Its contribution is the exact anchored canonical formulation: explicit finite-$n$ risk formulas under both the i.i.d.\ and fixed-composition models, and the integration of these formulas with the shuffle-privacy envelope theory developed in earlier sections.

\subsection{The ratio-capped template polytope}

\begin{proposition}[The raw local cap in template coordinates]\label{prop:raw-cap}
Fix $\lambda=e^{\varepsilon_0}\ge 1$. Let $\rho\in\fM_d$ and let $\nu=a_\#\rho$ be its template law. Then the channel $W_\rho$ is $\varepsilon_0$-LDP if and only if
\[
\nu(\cK_{d,\lambda})=1,
\qquad
\cK_{d,\lambda}:=\{a\in\cK_d:\ a_i\le \lambda a_j\ \text{for all }i,j\in[d]\}.
\]
\end{proposition}

\begin{proof}
For a template $a\in\cK_d$, the row-$i$ and row-$j$ output masses are $a_i$ and $a_j$ times the common average density. Thus
\[
\frac{dW_\rho(\cdot\mid j)}{dW_\rho(\cdot\mid i)}= \frac{a_j}{a_i}
\]
at every output point carrying template $a$. Therefore $W_\rho$ is $\varepsilon_0$-LDP if and only if
\[
e^{-\varepsilon_0}\le \frac{a_j}{a_i}\le e^{\varepsilon_0}
\qquad\text{for every }i,j
\]
on the support of $\nu$. Since $\lambda=e^{\varepsilon_0}$, this is equivalent to
\[
a_i\le \lambda a_j \qquad\forall\, i,j,
\]
that is, to $\nu(\cK_{d,\lambda})=1$.
\end{proof}

For $1\le s\le d-1$, define
\[
\beta_s(\lambda):=\frac{d}{d+s(\lambda-1)},
\qquad
\alpha_s(\lambda):=\lambda \beta_s(\lambda),
\]
and let $a^{(s)}_\lambda\in\cK_d$ be the sorted two-level template with $s$ coordinates equal to $\alpha_s(\lambda)$ and $d-s$ coordinates equal to $\beta_s(\lambda)$. The associated exchangeable orbit law is denoted by $\rho_{s,\lambda}$ and will be called the \emph{$s$-subset mechanism}. This mechanism class was introduced by Wang et al.\ \cite{wang-ldp}; see also \cite{ye-barg,pan-strict} for asymptotic and strict optimality results in the pure local model.

\begin{lemma}[Extreme templates under the ratio cap]\label{lem:raw-extreme}
Let $\lambda>1$. The extreme points of $\cK_{d,\lambda}$ are exactly the templates obtained, up to permutation, from $a^{(s)}_\lambda$ for some $s\in\{1,\dots,d-1\}$.
\end{lemma}

\begin{proof}
Let $a\in\cK_{d,\lambda}$, and write
\[
M:=\max_i a_i,
\qquad
m:=\min_i a_i.
\]
Because $a\in\cK_{d,\lambda}$, one has $0<m\le M\le \lambda m$.

We first show that if $M<\lambda m$, then $a$ is not extreme. Choose indices $p,q$ such that $a_p=M$ and $a_q=m$. Because the inequality is strict, there exists $\varepsilon>0$ such that
\[
m-\varepsilon>0,
\qquad
M+\varepsilon<\lambda(m-\varepsilon).
\]
Define
\[
a^\pm:=a\pm \varepsilon(e_p-e_q).
\]
Then $\sum_i a_i^\pm=d$ and all coordinates remain positive. All coordinates other than $p$ and $q$ are unchanged, and they satisfy the ratio constraints strictly because $M<\lambda m$ implies $a_i/a_j<\lambda$ for all $i,j$ with the same strict slack. Since the perturbed coordinates converge back to $(M,m)$ as $\varepsilon\downarrow 0$, the same strict slack persists for all sufficiently small $\varepsilon>0$. Hence $a^\pm\in\cK_{d,\lambda}$ and
\[
a=\frac12(a^++a^-)
\]
with $a^+\neq a^-$. Thus $a$ is not extreme. Therefore every extreme point must satisfy
\[
M=\lambda m.
\]

Now suppose $M=\lambda m$ but there exists an index $r$ with
\[
m<a_r<M.
\]
Let
\[
H:=\{i:\ a_i=M\},
\qquad
L:=\{i:\ a_i=m\},
\]
and set $s:=|H|$, $t:=|L|$. Define a direction $v\in\R^d$ by
\[
v_i=
\begin{cases}
\lambda, & i\in H,\\
1, & i\in L,\\
-(s\lambda+t), & i=r,\\
0, & \text{otherwise}.
\end{cases}
\]
Then $\sum_i v_i=0$. For every active max-min constraint $a_i=\lambda a_j$ with $i\in H$ and $j\in L$,
\[
v_i-\lambda v_j=\lambda-\lambda\cdot 1=0.
\]
All remaining ratio constraints are strict, because every coordinate outside $H\cup L$ lies strictly between $m$ and $M$. Since the perturbation direction $v$ preserves the active equalities to first order and all other constraints hold with strict slack at $a$, continuity implies that for all sufficiently small $\varepsilon>0$, both
\[
a^\pm:=a\pm \varepsilon v
\]
remain in $\cK_{d,\lambda}$, and
\[
a=\frac12(a^++a^-)
\]
with $a^+\neq a^-$. Hence $a$ is not extreme.

We have shown that every extreme point has exactly two distinct coordinate values, say $M$ and $m$, and that they satisfy $M=\lambda m$. Let $s$ be the number of coordinates equal to $M$. The sum constraint gives
\[
sM+(d-s)m=d,
\]
hence
\[
m=\frac{d}{d+s(\lambda-1)}=\beta_s(\lambda),
\qquad
M=\lambda m=\alpha_s(\lambda).
\]
Therefore every extreme point is, up to permutation, one of the templates $a^{(s)}_\lambda$.

Conversely, let $a$ be such a two-level template, with high set $H$ of size $s$ and low set $L$ of size $d-s$. Suppose
\[
a=t b+(1-t)c
\qquad\text{with }0<t<1,\quad b,c\in \cK_{d,\lambda}.
\]
Fix $i\in H$ and $j\in L$. Then
\[
a_i=M=\lambda m=\lambda a_j.
\]
Since $b$ and $c$ are feasible,
\[
b_i\le \lambda b_j,
\qquad
c_i\le \lambda c_j.
\]
Therefore
\[
M=a_i=t b_i+(1-t)c_i\le \lambda\bigl(t b_j+(1-t)c_j\bigr)=\lambda a_j=M.
\]
All inequalities are equalities, so
\[
b_i=\lambda b_j,
\qquad
c_i=\lambda c_j
\]
for every $i\in H$ and $j\in L$. Hence all high coordinates of $b$ are equal to some $\lambda \beta$, all low coordinates to $\beta$, and similarly for $c$. The sum constraint forces
\[
\beta=\frac{d}{d+s(\lambda-1)}=m,
\]
so $b=a$. The same argument gives $c=a$. Thus $a$ is extreme.
\end{proof}


\begin{theorem}[Exact raw-cap trace maximization]\label{thm:raw-trace}
Let $\lambda=e^{\varepsilon_0}>1$, and for $1\le s\le d-1$ define
\[
T_{d,\lambda}(s):=\frac{d\,s(d-s)(\lambda-1)^2}{(d+s(\lambda-1))^2},
\qquad
T_d^*(\lambda):=\max_{1\le s\le d-1} T_{d,\lambda}(s).
\]
Then
\[
\sup_{\rho:\,W_\rho\ \varepsilon_0\text{-LDP}}\tr\Sigma(\rho)=T_d^*(\lambda).
\]
Let
\[
\cS_{d,\lambda}^*:=\argmax_{1\le s\le d-1} T_{d,\lambda}(s).
\]
For every $s_*\in \cS_{d,\lambda}^*$, the exchangeable $s_*$-subset mechanism $\rho_{s_*,\lambda}$ is optimal. More generally, an exchangeable law is optimal if and only if it is a convex combination of the laws $\{\rho_{s,\lambda}:s\in \cS_{d,\lambda}^*\}$; in particular, the extreme exchangeable optimizers are exactly the individual maximizing subset mechanisms. Finally,
\[
\cS_{d,\lambda}^*\subseteq \left\{\left\lfloor \frac{d}{\lambda+1}\right\rfloor,\left\lceil \frac{d}{\lambda+1}\right\rceil\right\}\cap \{1,\dots,d-1\}.
\]
\end{theorem}

\begin{proof}
Let $\rho$ be feasible. By \cref{prop:raw-cap}, its template law is supported on $\cK_{d,\lambda}$. Symmetrizing does not change $\tr\Sigma$ and preserves the support constraint, so by replacing $\rho$ with $\bar\rho$ we may assume that $\rho$ is exchangeable. Let $\nu$ be the corresponding law on the sorted chamber. Then
\[
\tr\Sigma(\rho)=\int (B(a)-d)\,\nu(da)\le \max_{a\in \cK_{d,\lambda}} (B(a)-d),
\]
with equality if and only if $\nu$ is supported on the set of maximizers of $B(a)-d$ over $\cK_{d,\lambda}$.

Since $B(a)-d$ is strictly convex on the convex polytope $\cK_{d,\lambda}$, every maximizer is an extreme point. By \cref{lem:raw-extreme}, each extreme point is one of the templates $a^{(s)}_\lambda$. For $a^{(s)}_\lambda$, a direct computation gives
\[
B(a^{(s)}_\lambda)-d
=
s\alpha_s(\lambda)^2+(d-s)\beta_s(\lambda)^2-d
=
\frac{d\,s(d-s)(\lambda-1)^2}{(d+s(\lambda-1))^2}
=
T_{d,\lambda}(s).
\]
Therefore
\[
\max_{a\in \cK_{d,\lambda}}(B(a)-d)=T_d^*(\lambda),
\]
which proves the exact formula for the optimum. The support characterization above shows that an exchangeable law is optimal exactly when its sorted-template law is supported on
\[
\{a^{(s)}_\lambda:s\in \cS_{d,\lambda}^*\},
\]
equivalently when it is a convex combination of the orbit laws $\rho_{s,\lambda}$ with $s\in \cS_{d,\lambda}^*$. The extreme exchangeable optimizers are therefore precisely the individual maximizing subset mechanisms.

It remains to locate the maximizing integer set. Consider the continuous function
\[
g_{d,\lambda}(s):=\frac{s(d-s)}{(d+s(\lambda-1))^2},
\qquad s\in [1,d-1].
\]
Since $T_{d,\lambda}(s)=d(\lambda-1)^2 g_{d,\lambda}(s)$, maximizing $T_{d,\lambda}$ is equivalent to maximizing $g_{d,\lambda}$. Differentiating,
\[
g'_{d,\lambda}(s)=\frac{d(d-(\lambda+1)s)}{(d+s(\lambda-1))^3}.
\]
Thus $g_{d,\lambda}$ is strictly increasing on
\[
\left[1,\frac{d}{\lambda+1}\right]
\]
and strictly decreasing on
\[
\left[\frac{d}{\lambda+1},d-1\right].
\]
Hence every maximizing integer lies among the two adjacent integers bracketing $d/(\lambda+1)$, namely
\[
\cS_{d,\lambda}^*\subseteq \left\{\left\lfloor \frac{d}{\lambda+1}\right\rfloor,\left\lceil \frac{d}{\lambda+1}\right\rceil\right\}\cap \{1,\dots,d-1\}.
\]
\end{proof}

\begin{theorem}[Exact canonical design under the raw local cap]\label{thm:raw-canonical}
Let $\lambda=e^{\varepsilon_0}>1$ and let $T_d^*(\lambda)$ be as in \cref{thm:raw-trace}. Then
\[
\inf_{\rho:\,W_\rho\ \varepsilon_0\text{-LDP}}
\sup_{\theta\in\Delta_d}
R^{\mathrm{iid}}_\theta(\rho)
=
\frac{(d-1)^2}{n\,T_d^*(\lambda)},
\]
and
\[
\inf_{\rho:\,W_\rho\ \varepsilon_0\text{-LDP}}
\sup_{\theta\in\Delta_d}
R^{\mathrm{fc}}_\theta(\rho)
=
\frac1n\left[
\frac{(d-1)^2}{T_d^*(\lambda)}
-\frac{d-1}{d}
\right].
\]
Let $\cS_{d,\lambda}^*$ be the maximizing set from \cref{thm:raw-trace}. Both values are attained by every exchangeable law whose sorted-template law is supported on $\{a^{(s)}_\lambda:s\in \cS_{d,\lambda}^*\}$; in particular every maximizing $s_*$-subset mechanism is optimal. Conversely, every extreme exchangeable optimizer is one of these subset-selection laws, and every exchangeable optimizer is their convex combination.
\end{theorem}

\begin{proof}
By \cref{prop:raw-cap}, the raw local cap is preserved under symmetrization. Therefore \cref{thm:symmetrization,cor:iid-symm} allow us to restrict both optimization problems to exchangeable laws.

If $\rho$ is exchangeable, then \cref{thm:symmetrization}(iii) gives
\[
\Sigma(\rho)=\frac{\tr\Sigma(\rho)}{d-1}\,I_{d-1}.
\]
Hence \cref{prop:iid-risk} yields
\[
\sup_{\theta\in\Delta_d} R^{\mathrm{iid}}_\theta(\rho)
=
\frac{(d-1)^2}{n\,\tr\Sigma(\rho)},
\]
while \cref{prop:fixed-risk} gives
\[
\sup_{\theta\in\Delta_d} R^{\mathrm{fc}}_\theta(\rho)
=
\frac1n\left[
\frac{(d-1)^2}{\tr\Sigma(\rho)}
-\frac{d-1}{d}
\right].
\]
Both quantities are strictly decreasing functions of $\tr\Sigma(\rho)$, so minimizing the risks is equivalent to maximizing $\tr\Sigma(\rho)$ under the raw cap. Theorem~\ref{thm:raw-trace} identifies the maximum as $T_d^*(\lambda)$ and characterizes exactly which exchangeable laws attain it. This proves the claim.
\end{proof}

\begin{remark}[Scope of the optimizer characterization]
\cref{thm:raw-trace,thm:raw-canonical} characterize exchangeable optimizers. They do not classify all raw-cap optimizers; \cref{thm:symmetrization} only shows that an exchangeable optimizer always exists and suffices for the canonical risk problem.
\end{remark}

\begin{corollary}[Asymptotic minimax under a varying raw cap]\label{cor:raw-asymptotic}
Let $\lambda_n=e^{\varepsilon_{0,n}}>1$ and assume
\[
n\,T_d^*(\lambda_n)\longrightarrow\infty.
\]
Then
\[
\inf_{\rho:\,W_\rho\ \varepsilon_{0,n}\text{-LDP}}
\inf_{\widehat\theta}
\sup_{\theta\in\Delta_d}
\E_{\theta,\rho}\|\widehat\theta-\theta\|_2^2
=
\frac{(d-1)^2}{n\,T_d^*(\lambda_n)}\,(1+o(1)).
\]
\end{corollary}

\begin{proof}
The upper bound is provided by \cref{thm:raw-canonical}. For the lower bound, repeat the proof of \cref{thm:main-design} with the trace-cap estimate replaced by the exact raw-cap estimate
\[
\tr\Sigma(\rho)\le T_d^*(\lambda_n),
\]
which follows from \cref{thm:raw-trace}. The same local prior with
\[
r_n:=\bigl(n\,T_d^*(\lambda_n)\bigr)^{-1/4}
\]
satisfies $r_n\to 0$ by assumption, and the remainder of the van Trees argument is unchanged.
\end{proof}

\subsection{Numerical illustration}

For the grid $d\in \{3,5,10,20\}$ and $\varepsilon_0\in \{0.5,1,2\}$, Table~\ref{tab:raw-phase} in Appendix~\ref{app:tables} reports the maximizing subset size, the exact raw-cap trace optimum $T_d^*(e^{\varepsilon_0})$, and the exact finite-$n$ canonical risk constants from \cref{thm:raw-canonical}. The actual i.i.d.\ and fixed-composition risks are the tabulated constants divided by $n$. Table~\ref{tab:raw-cap} compares the exact raw-cap trace maximum with the low-budget trace cap $K_d\,\chi_*(\rho)$ from Section~\ref{sec:tracecap} and a crude bound. The exact raw-cap trace lies noticeably below the low-budget $\chi_*$ cap once $\varepsilon_0$ is no longer small, whereas the crude bound is uniformly far too large.



\section{Discussion}\label{sec:discussion}

The anchored affine law now gives a complete exact language for three tightly connected questions:
\begin{enumerate}[label=\textup{(\roman*)}]
\item exact pairwise privacy envelopes under shuffling, together with a rigidity converse;
\item exact canonical design under the canonical pairwise $\chi_*$ budget, both at finite $n$ and in the low-budget asymptotic minimax regime;
\item exact canonical design under the raw local $\varepsilon_0$ cap, with an explicit subset-selection phase diagram.
\end{enumerate}

Within this scope the theory is closed. The anchored law $\rho$ is the exact channel carrier, projective fibers reduce pairwise privacy to scalar shadows, the covariance matrix $\Sigma(\rho)$ governs the mechanism-calibrated canonical estimator, and the template orbit curve $\Gamma_d$ governs the full $\chi_*$ frontier. The paper also resolves two structural questions left open in earlier drafts: the exact finite-$n$ canonical optimum under the $\chi_*$ budget is explicit, and the global $\chi_*$ frontier is always realized by a mixture of at most two exchangeable orbit laws. In particular, the two-orbit theorem closes the global $\chi_*$-frontier problem that remained open in \cite{shvets-growing} beyond the low-budget tangency point.

The raw local cap problem is now equally explicit at the canonical level. \emph{Augmented randomized response is not the raw-cap optimizer.} The correct optimizer is an exchangeable subset-selection law, with maximizing subset size determined by the phase set
\[
\cS_{d,\lambda}^*=\argmax_{1\le s\le d-1}\frac{d\,s(d-s)(\lambda-1)^2}{(d+s(\lambda-1))^2},
\qquad \lambda=e^{\varepsilon_0}.
\]
Thus the raw-cap and $\chi_*$-budget design problems are genuinely different, even though both become exact moment problems in the anchored gauge. The subset-selection mechanism class was introduced by Wang, Blocki, Li, and Jha \cite{wang-ldp}, and Ye and Barg \cite{ye-barg} showed that the optimal subset size is asymptotically one of the nearest integers to $d/(e^{\varepsilon_0}+1)$. The present contribution is the exact anchored canonical formulation---explicit finite-$n$ risk formulas, the trace-cap integration, and the full phase-diagram packaging---not the first identification of the optimizer class.

A separate comparison is with Takagi and Liew \cite{takagi-liew}. The two programs address overlapping conceptual territory---both seek low-dimensional mechanism summaries that control shuffled privacy---but differ in their main objects and results. Takagi--Liew work with blanket divergence, extract a one-parameter shuffle index, and obtain asymptotic upper and lower bands beyond pure local differential privacy, together with an FFT-based finite-$n$ accountant. The present paper works with anchored laws, exact finite-$n$ universal envelopes for all convex $f$-divergences, rigidity converses, and canonical design frontiers. Conversely, the present paper does not attempt to recover the broader non-pure-LDP asymptotic regime of \cite{takagi-liew} from the anchored gauge. The two programs address related privacy-summary questions at a different level of generality and with different technical objects.

Concurrently and independently, Pan \cite{pan-strict} proved strict optimality of subset-selection mechanisms for frequency estimation under local differential privacy from a communication-efficiency perspective, building on the Wang--Ye--Barg line \cite{wang-ldp,ye-barg}. Pan's strict optimality result in the pure local model directly implies that the optimizer class identified in Section~\ref{sec:raw} is not a new discovery; the present contribution is the exact anchored canonical formulation---explicit finite-$n$ risk formulas, the trace-cap integration, and the connection with the shuffle-privacy envelope theory developed in earlier sections.

The main limitation remains asymptotic. The anchored law and its projective fibers do \emph{not} subsume the author's fixed-composition non-Gaussian asymptotics from \cite{shvets-part2,shvets-part3}. For interior base points $h\in \operatorname{int}\cH$, \cref{prop:fiber-props} shows that the fiber $\eta_h$ is bounded. Therefore any asymptotic theory built only from the i.i.d.\ product family $q_h^{\otimes n}$ and the scaled fibers $Y_h/r_n$ is necessarily Gaussian in the interior. The Skellam and hybrid regimes of the trilogy must therefore live on a different carrier: the fixed-composition quotient geometry and its conditional likelihood field.

In summary, anchored laws are the right exact coordinate system for privacy and canonical design; the trilogy's quotient geometry remains the right coordinate system for the non-Gaussian fixed-composition asymptotics.

Several questions remain open.
\begin{enumerate}[label=\textup{(\arabic*)}]
\item Give an explicit closed form for the upper concave envelope of $\Gamma_d$ in the high-budget $\chi_*$ regime.
\item Extend the raw-cap and $\chi_*$ symmetrization arguments to fixed-message multi-message shuffle protocols.
\item Determine whether the exact finite-$n$ global minimax constant under the $\chi_*$ budget, over all estimators, coincides with the canonical finite-$n$ formula proved here.
\item Develop the asymptotic unification of Parts~I--III in the correct quotient geometry, rather than in the anchored gauge.
\end{enumerate}


\appendix


\section{A multivariate van Trees bound}\label{app:vantrees}

\begin{lemma}[Coordinatewise van Trees]\label{lem:vantrees}
Let $U=(U_1,\dots,U_p)$ have prior density $\pi$ on a rectangle $\prod_{j=1}^p[a_j,b_j]$, where $\pi$ is continuously differentiable, strictly positive in the interior, and vanishes on the boundary faces. Let $Y$ be an observation with conditional density $p_u(y)$, and let $\widehat U(Y)$ be any estimator. Assume the score functions are square-integrable and differentiation under the integral sign is justified. Then
\[
\int \E_u\|\widehat U-U\|_2^2\,\pi(u)\,du
\ge
\frac{p^2}{\int \tr I(u)\,\pi(u)\,du+J(\pi)},
\]
where
\[
I(u)=\E_u\bigl[(\nabla_u\log p_u(Y))(\nabla_u\log p_u(Y))^\top\bigr],
\qquad
J(\pi)=\int \frac{\|\nabla \pi(u)\|_2^2}{\pi(u)}\,du.
\]
\end{lemma}

\begin{proof}
Fix a coordinate $j\in\{1,\dots,p\}$. Under the joint law $\pi(u)p_u(y)\,du\,dy$,
\[
\E\Bigl[(\widehat U_j-U_j)\bigl(\partial_j\log p_U(Y)+\partial_j\log \pi(U)\bigr)\Bigr]=1.
\]
Indeed,
\[
\begin{aligned}
&\int \!\!\int (\widehat u_j(y)-u_j)\bigl(\partial_j p_u(y)\,\pi(u)+p_u(y)\partial_j\pi(u)\bigr)\,dy\,du\\
&\qquad = -\int \!\!\int \partial_j(\widehat u_j(y)-u_j)\,p_u(y)\pi(u)\,dy\,du
=\int\!\!\int p_u(y)\pi(u)\,dy\,du=1,
\end{aligned}
\]
where integration by parts is justified by the boundary condition $\pi=0$ on the boundary faces.

By Cauchy--Schwarz,
\[
\E[(\widehat U_j-U_j)^2]\ge \frac{1}{\E[(S_j+Q_j)^2]},
\]
where
\[
S_j:=\partial_j\log p_U(Y),\qquad Q_j:=\partial_j\log \pi(U).
\]
Since $\E[S_j\mid U]=0$, the cross term vanishes and
\[
\E[(S_j+Q_j)^2]=\int I_{jj}(u)\,\pi(u)\,du+\int \frac{(\partial_j\pi(u))^2}{\pi(u)}\,du.
\]
Summing over $j$ and using the harmonic-arithmetic mean inequality,
\[
\sum_{j=1}^p \E[(\widehat U_j-U_j)^2]
\ge \sum_{j=1}^p \frac{1}{A_j+B_j}
\ge \frac{p^2}{\sum_j(A_j+B_j)},
\]
where
\[
A_j:=\int I_{jj}(u)\,\pi(u)\,du,
\qquad
B_j:=\int \frac{(\partial_j\pi(u))^2}{\pi(u)}\,du.
\]
Since $\sum_j A_j=\int \tr I(u)\,\pi(u)\,du$ and $\sum_j B_j=J(\pi)$, the result follows.
\end{proof}


\section{Binary randomized response and the endpoint law}\label{app:brr}

Binary randomized response with local parameter $\varepsilon_0$ is the two-row channel
\[
W_{\mathrm{BRR}}=
\begin{pmatrix}
\dfrac{e^{\varepsilon_0}}{1+e^{\varepsilon_0}} & \dfrac{1}{1+e^{\varepsilon_0}}\\[1.2ex]
\dfrac{1}{1+e^{\varepsilon_0}} & \dfrac{e^{\varepsilon_0}}{1+e^{\varepsilon_0}}
\end{pmatrix}.
\]
Under row $1$, the likelihood ratio of row $2$ against row $1$ takes the values $e^{-\varepsilon_0}$ and $e^{\varepsilon_0}$ with exactly the weights defining the endpoint law $\mu^\star$ from Section~\ref{sec:privacy}. This is the only fact from the binary case used in the proof of the universal envelope.

\section{Extreme points of moment sets and of \texorpdfstring{$\fM_d$}{Md}}\label{app:extreme}

The support bound below is classical; see Winkler \cite{winkler-moment} for general extreme-point results on moment sets. We include a short proof specialized to our setting for completeness.

\begin{theorem}[Support theorem for linear moment sets]\label{thm:moment-support}
Let $K$ be a separable metrizable space equipped with its Borel $\sigma$-field, let $f_1,\dots,f_m:K\to \R$ be measurable, and fix $c=(c_1,\dots,c_m)\in \R^m$. Define
\[
\cP_f(c):=\Bigl\{\nu\in \cP(K):\ \int |f_r|\,d\nu<\infty,\ \int f_r\,d\nu=c_r\ \text{for }r=1,\dots,m\Bigr\}.
\]
Then every extreme point of $\cP_f(c)$ is supported on at most $m+1$ points.

Conversely, if
\[
\nu=\sum_{k=1}^N p_k\,\delta_{x_k},
\qquad
1\le N\le m+1,
\qquad
p_k>0,
\]
and the vectors
\[
\bigl(1,f_1(x_k),\dots,f_m(x_k)\bigr)\in \R^{m+1},
\qquad k=1,\dots,N,
\]
are linearly independent, then $\nu$ is an extreme point of $\cP_f(c)$.
\end{theorem}

\begin{proof}
Suppose first that $\nu\in \cP_f(c)$ is extreme and that its support contains at least $m+2$ distinct points $x_1,\dots,x_{m+2}$. Because $K$ is metrizable, we may choose pairwise disjoint open neighborhoods $E_1,\dots,E_{m+2}$ of these points such that
\[
\nu(E_k)>0
\qquad (k=1,\dots,m+2).
\]
For each $k$, define the vector
\[
v_k:=\left(\nu(E_k),\ \int_{E_k} f_1\,d\nu,\ \dots,\ \int_{E_k} f_m\,d\nu\right)\in \R^{m+1}.
\]
Since there are $m+2$ vectors in $\R^{m+1}$, they are linearly dependent. Thus there exist real numbers $\alpha_1,\dots,\alpha_{m+2}$, not all zero, such that
\[
\sum_{k=1}^{m+2}\alpha_k v_k=0.
\]
Equivalently,
\[
\sum_{k=1}^{m+2}\alpha_k \nu(E_k)=0,
\qquad
\sum_{k=1}^{m+2}\alpha_k \int_{E_k} f_r\,d\nu=0
\quad (r=1,\dots,m).
\]

Define the signed measure
\[
\sigma:=\sum_{k=1}^{m+2}\alpha_k\,\nu|_{E_k}.
\]
Then
\[
\sigma(K)=0,
\qquad
\int f_r\,d\sigma=0
\quad (r=1,\dots,m).
\]
Choose
\[
\varepsilon:=\min_{\alpha_k\neq 0}\frac{1}{2|\alpha_k|}>0.
\]
Since the sets $E_k$ are disjoint, the measures
\[
\nu_\pm:=\nu\pm \varepsilon \sigma
\]
are both nonnegative. Indeed, on each $E_k$ the Radon--Nikodym factor is $1\pm \varepsilon\alpha_k\ge 1/2$, and outside $\bigcup_k E_k$ the measure is unchanged. Moreover,
\[
\nu_\pm(K)=1,
\qquad
\int f_r\,d\nu_\pm=c_r
\quad (r=1,\dots,m),
\]
so $\nu_\pm\in \cP_f(c)$. Since $\sigma\neq 0$, one has $\nu_+\neq \nu_-$, and
\[
\nu=\frac12(\nu_++\nu_-),
\]
contradicting extremality. Hence every extreme point has support size at most $m+1$.

Conversely, let
\[
\nu=\sum_{k=1}^N p_k\,\delta_{x_k}
\]
with linearly independent moment vectors as in the statement, and suppose
\[
\nu=t\nu_1+(1-t)\nu_2,
\qquad
0<t<1,
\qquad
\nu_1,\nu_2\in \cP_f(c).
\]
Since $\nu$ is supported on $\{x_1,\dots,x_N\}$, both $\nu_1$ and $\nu_2$ are also supported on this set: if $B$ is disjoint from $\{x_1,\dots,x_N\}$, then
\[
0=\nu(B)=t\nu_1(B)+(1-t)\nu_2(B)
\]
forces $\nu_1(B)=\nu_2(B)=0$.

Write
\[
\nu_\ell=\sum_{k=1}^N q_k^{(\ell)}\,\delta_{x_k},
\qquad \ell=1,2.
\]
The constraints $\nu_\ell\in \cP_f(c)$ are exactly
\[
\sum_{k=1}^N q_k^{(\ell)}=1,
\qquad
\sum_{k=1}^N q_k^{(\ell)} f_r(x_k)=c_r
\quad (r=1,\dots,m).
\]
Since the vectors $(1,f_1(x_k),\dots,f_m(x_k))$ are linearly independent, this linear system has a unique solution. But the weights $p_k$ solve it, so
\[
q_k^{(1)}=q_k^{(2)}=p_k
\qquad (k=1,\dots,N).
\]
Thus $\nu_1=\nu_2=\nu$, proving that $\nu$ is extreme.
\end{proof}

\begin{corollary}[Extreme points of $\fM_d$]\label{cor:extreme-Md}
A measure $\rho\in \fM_d$ is an extreme point of $\fM_d$ if and only if it is finitely supported on an affinely independent subset of $\cX_d$. In particular, every extreme point of $\fM_d$ has support size at most $d$.
\end{corollary}

\begin{proof}
Apply \cref{thm:moment-support} with $m=d-1$, $K=\cX_d$, and
\[
f_r(x)=x_r,\qquad r=1,\dots,d-1,
\]
together with the moment condition
\[
\int x\,\rho(dx)=0.
\]
The support bound becomes $d$ points. For a finitely supported measure
\[
\rho=\sum_{k=1}^N p_k\,\delta_{x_k},\qquad p_k>0,
\]
the vectors
\[
(1,f_1(x_k),\dots,f_{d-1}(x_k))=(1,x_k)\in \R^d
\]
are linearly independent exactly when the points $x_1,\dots,x_N$ are affinely independent. The converse criterion in \cref{thm:moment-support} therefore shows that affine independence implies extremality.

For necessity, suppose $\rho=\sum_{k=1}^N p_k\,\delta_{x_k}\in\fM_d$ with $p_k>0$ and the points $x_1,\dots,x_N$ are affinely dependent. Then there exist real numbers $c_1,\dots,c_N$, not all zero, such that
\[
\sum_{k=1}^N c_k=0,\qquad \sum_{k=1}^N c_k x_k=0.
\]
For all sufficiently small $\varepsilon>0$, the measures
\[
\rho_\pm:=\sum_{k=1}^N (p_k\pm\varepsilon c_k)\,\delta_{x_k}
\]
are distinct probability measures in $\fM_d$, and $\rho=\tfrac12(\rho_++\rho_-)$. Hence $\rho$ is not extreme. Therefore every extreme point of $\fM_d$ is supported on an affinely independent set.
\end{proof}


\section{Numerical tables}\label{app:tables}

Table~\ref{tab:raw-phase} reports the exact subset-selection phase diagram and canonical risk constants for a grid of alphabet sizes and local privacy levels. Table~\ref{tab:raw-cap} compares the exact raw-cap trace maximum with the low-budget trace cap
\[
K_d\,\chi_*(\rho),\qquad K_d:=\frac{d(d-1)}{d+2\sqrt{d-1}},
\]
and the crude bound from the trivial pointwise inequality $(B(a)-d)/(A(a)B(a)-d^2)\le 1$.

\begin{table}[ht]
\centering
\small
\caption{Exact subset-selection phase diagram and canonical risk constants. The listed values are the dimensionless products $nR_{\max}^{\mathrm{iid}}$ and $nR_{\max}^{\mathrm{fc}}$; the actual risks equal these constants divided by $n$. Numerical values are rounded to four decimal places.}
\label{tab:raw-phase}
\begin{tabular}{@{}cccccc@{}}
\toprule
$d$ & $\varepsilon_0$ & $\arg\max s$ & $T_d^*(e^{\varepsilon_0})$ & $nR_{\max}^{\mathrm{iid}}$ & $nR_{\max}^{\mathrm{fc}}$ \\
\midrule
3 & 0.5 & 1 & 0.1897 & 21.0899 & 20.4232 \\
3 & 1 & 1 & 0.7957 & 5.0268 & 4.3601 \\
3 & 2 & 1 & 2.7783 & 1.4397 & 0.7731 \\
\addlinespace
5 & 0.5 & 2 & 0.3184 & 50.2587 & 49.4587 \\
5 & 1 & 1 & 1.3083 & 12.2298 & 11.4298 \\
5 & 2 & 1 & 6.2940 & 2.5421 & 1.7421 \\
\addlinespace
10 & 0.5 & 4 & 0.6367 & 127.2172 & 126.3172 \\
10 & 1 & 3 & 2.6996 & 30.0041 & 29.1041 \\
10 & 2 & 1 & 13.6775 & 5.9221 & 5.0221 \\
\addlinespace
20 & 0.5 & 8 & 1.2734 & 283.4902 & 282.5402 \\
20 & 1 & 5 & 5.4176 & 66.6344 & 65.6844 \\
20 & 2 & 2 & 27.3551 & 13.1968 & 12.2468 \\
\bottomrule
\end{tabular}
\end{table}

\begin{table}[ht]
\centering
\small
\caption{Exact raw-cap trace maximum versus the low-budget trace cap and a crude bound, all evaluated at a maximizing subset-selection law. Numerical values are rounded to four decimal places.}
\label{tab:raw-cap}
\begin{tabular}{@{}cccccc@{}}
\toprule
$d$ & $\varepsilon_0$ & $\arg\max s$ & exact $T_d^*(e^{\varepsilon_0})$ & $K_d\chi_*$ & $d(d-1)\chi_*$ \\
\midrule
3 & 0.5 & 1 & 0.1897 & 0.1907 & 1.1118 \\
3 & 1 & 1 & 0.7957 & 0.8812 & 5.1358 \\
3 & 2 & 1 & 2.7783 & 5.0813 & 29.6160 \\
\addlinespace
5 & 0.5 & 2 & 0.3184 & 0.3579 & 3.2208 \\
5 & 1 & 1 & 1.3083 & 1.3359 & 12.0229 \\
5 & 2 & 1 & 6.2940 & 9.0427 & 81.3841 \\
\addlinespace
10 & 0.5 & 4 & 0.6367 & 0.8052 & 12.8832 \\
10 & 1 & 3 & 2.6996 & 3.4977 & 55.9634 \\
10 & 2 & 1 & 13.6775 & 15.9062 & 254.4990 \\
\addlinespace
20 & 0.5 & 8 & 1.2734 & 1.7944 & 51.5326 \\
20 & 1 & 5 & 5.4176 & 7.3780 & 211.8811 \\
20 & 2 & 2 & 27.3551 & 35.4483 & 1017.9961 \\
\bottomrule
\end{tabular}
\end{table}


\begin{thebibliography}{99}

\bibitem{shvets-part1}
A.~Shvets.
\newblock Universal shuffle asymptotics: Sharp privacy analysis in the Gaussian regime.
\newblock \texttt{arXiv:2602.09029}, 2026.

\bibitem{shvets-part2}
A.~Shvets.
\newblock Universal shuffle asymptotics, Part~II: Non-Gaussian limits for shuffle privacy --- Poisson, Skellam, and compound-Poisson regimes.
\newblock \texttt{arXiv:2603.10073}, 2026.

\bibitem{shvets-part3}
A.~Shvets.
\newblock Universal shuffle asymptotics, Part~III: Dominant-block quotient geometry and hybrid Gaussian--compound-Poisson limits in finite-alphabet shuffle privacy.
\newblock \texttt{arXiv:2603.13407}, 2026.

\bibitem{shvets-growing}
A.~Shvets.
\newblock Growing alphabets do not automatically amplify shuffle privacy: obstruction, estimation bounds, and optimal mechanism design.
\newblock \texttt{arXiv:2603.18080}, 2026.

\bibitem{balle-blanket}
B.~Balle, J.~Bell, A.~Gasc\'on, and K.~Nissim.
\newblock The privacy blanket of the shuffle model.
\newblock In \emph{Advances in Cryptology --- CRYPTO 2019}, volume 11693 of \emph{Lecture Notes in Computer Science}, pages 638--667. Springer, 2019.

\bibitem{feldman-clones}
V.~Feldman, A.~McMillan, and K.~Talwar.
\newblock Hiding among the clones: A simple and nearly optimal analysis of privacy amplification by shuffling.
\newblock In \emph{Proceedings of the 62nd IEEE Symposium on Foundations of Computer Science}, pages 954--964. IEEE, 2021.

\bibitem{feldman-stronger}
V.~Feldman, A.~McMillan, and K.~Talwar.
\newblock Stronger privacy amplification by shuffling for R\'enyi and approximate differential privacy.
\newblock In \emph{Proceedings of the 2023 Annual ACM-SIAM Symposium on Discrete Algorithms}, pages 4966--4981. SIAM, 2023.

\bibitem{girgis-rdp}
A.~M. Girgis, D.~Data, S.~Diggavi, A.~T. Suresh, and P.~Kairouz.
\newblock On the R\'enyi differential privacy of the shuffle model.
\newblock In \emph{Proceedings of the 2021 ACM SIGSAC Conference on Computer and Communications Security}, pages 2321--2341. ACM, 2021.

\bibitem{koskela-shuffle}
A.~Koskela, A.~Heikkil\"a, and A.~Honkela.
\newblock Numerical accounting in the shuffle model of differential privacy.
\newblock \emph{Transactions on Machine Learning Research}, 2023.

\bibitem{wang-unified}
S.~Wang, Y.~Peng, J.~Li, Z.~Wen, Z.~Li, S.~Yu, D.~Wang, and W.~Yang.
\newblock Privacy amplification via shuffling: Unified, simplified, and tightened.
\newblock \emph{Proceedings of the VLDB Endowment}, 17(8):1870--1883, 2024.

\bibitem{erlingsson-shuffle}
\'U.~Erlingsson, V.~Feldman, I.~Mironov, A.~Raghunathan, K.~Talwar, and A.~Thakurta.
\newblock Amplification by shuffling: From local to central differential privacy via anonymity.
\newblock In \emph{Proceedings of the Thirtieth Annual ACM-SIAM Symposium on Discrete Algorithms}, pages 2468--2479. SIAM, 2019.

\bibitem{cheu-shuffle}
A.~Cheu, A.~Smith, J.~Ullman, D.~Zeber, and M.~Zhilyaev.
\newblock Distributed differential privacy via shuffling.
\newblock In \emph{Advances in Cryptology --- EUROCRYPT 2019, Part~I}, volume 11476 of \emph{Lecture Notes in Computer Science}, pages 375--403. Springer, 2019.

\bibitem{takagi-liew}
S.~Takagi and S.~P. Liew.
\newblock Analysis of shuffling beyond pure local differential privacy.
\newblock \texttt{arXiv:2601.19154}, 2026. Accepted to PODS 2026.

\bibitem{pan-strict}
M.~Pan.
\newblock Strict optimality of frequency estimation under local differential privacy.
\newblock \texttt{arXiv:2603.11523}, 2026.

\bibitem{strassen}
V.~Strassen.
\newblock The existence of probability measures with given marginals.
\newblock \emph{Annals of Mathematical Statistics}, 36(2):423--439, 1965.

\bibitem{kov}
P.~Kairouz, S.~Oh, and P.~Viswanath.
\newblock Extremal mechanisms for local differential privacy.
\newblock \emph{Journal of Machine Learning Research}, 17(17):1--51, 2016.

\bibitem{wang-ldp}
T.~Wang, J.~Blocki, N.~Li, and S.~Jha.
\newblock Locally differentially private protocols for frequency estimation.
\newblock In \emph{Proceedings of the 26th USENIX Security Symposium}, pages 729--745. USENIX, 2017.

\bibitem{ye-barg}
M.~Ye and A.~Barg.
\newblock Optimal schemes for discrete distribution estimation under locally differential privacy.
\newblock \emph{IEEE Transactions on Information Theory}, 64(8):5662--5676, 2018.

\bibitem{vantrees}
H.~L. Van Trees.
\newblock \emph{Detection, Estimation, and Modulation Theory, Part~I}.
\newblock Wiley, 1968.

\bibitem{torgersen}
E.~Torgersen.
\newblock \emph{Comparison of Statistical Experiments}.
\newblock Cambridge University Press, 1991.

\bibitem{winkler-moment}
G.~Winkler.
\newblock Extreme points of moment sets.
\newblock \emph{Mathematics of Operations Research}, 13(4):581--587, 1988.

\bibitem{blackwell-comparison}
D.~Blackwell.
\newblock Equivalent comparisons of experiments.
\newblock \emph{Annals of Mathematical Statistics}, 24(2):265--272, 1953.

\end{thebibliography}

\end{document}
